\chapter{选品、供应链与履约}
\label{ch:supply}

\begin{insight}[本章要解决的问题]
\begin{enumerate}
  \item 跨境零售的胜负被概括为「七分在选品与供应链，三分在流量」。这句话如果成立，那么\emph{选品}到底是一个直觉行为，还是一个可以分解、可以打分、可以被证伪的决策流程？
  \item 一个候选商品有搜索量、有榜单排名、也有看起来不错的毛利率，为什么上线三个月后仍然亏钱？\emph{全链路单位经济}里通常被漏掉的是哪几项？
  \item 1688 上有几万家看起来一模一样的供应商，凭什么判断哪一家是工厂、哪一家是贸易商、哪一家会在旺季掉链子？\kw{打样}与\kw{验厂}应该看什么、花多少钱、值不值得？
  \item AQL 抽样里的 Ac/Re 是什么意思？供应商说「我们按 2.5 收货」，这句话在统计上到底约束了什么、又放过了什么？
  \item 安全库存公式 $SS = z\sqrt{LT\sigma_d^2 + d^2\sigma_{LT}^2}$ 里的四个量，哪一个最容易算错？为什么把服务水平从 95\% 提到 99\%，库存要涨四成而缺货只少四个百分点？
  \item 海外仓应该靠近港口还是靠近消费者？空运补货在什么条件下才比海运划算？「最后一公里」为什么能吃掉物流成本的一半？
  \item 牛鞭效应是供应链的天然属性还是管理失误的产物？当关税上调 25\% 时，正确的动作是涨价、换产地，还是改产品？
\end{enumerate}
\end{insight}

本章讨论跨境电商里最不性感、但决定生死的那一半：\kw{把东西做出来、运过去、卖出去、并且赚到钱}。
前面的章节讨论的是需求侧——搜索、推荐、广告、内容，它们的共同点是\emph{试错成本低、反馈快}：
一个广告素材跑不动，两小时就能关掉；一个关键词排不上去，换一个就是。
而供应链的动作是反过来的：\textbf{反馈慢（以周甚至月计）、单次决策金额大、且几乎不可逆}。
一个错误的备货决策，会在四个月后以「仓库里躺着一万件卖不动的货」的形式出现在财务报表上。

因此本章的方法论重心不是「找到爆款」，而是\kw{用闸门把不可逆决策切成可逆的小步}。
我们沿着一条真实的链路推进：选品的五道闸门（\cref{sec:sup-select}）、供应商的寻找与治理（\cref{sec:sup-supplier}）、
品控与合规（\cref{sec:sup-qc}）、库存与计划（\cref{sec:sup-inventory}）、
履约与仓配（\cref{sec:sup-fulfill}）、供应链韧性（\cref{sec:sup-resilience}），
最后是支撑这一切的数字化底座（\cref{sec:sup-digital}）。

本章的引用与数据约定：库存与补货理论以 \cite{silver1998inventory} 与 \cite{chopra2019supply} 的经典结论为准，
预测方法以 \cite{hyndman2021forecasting} 为准，牛鞭效应的原始论述来自 \cite{lee1997bullwhip}，
约束理论来自 \cite{goldratt2004goal}，国际贸易术语以 ICC 的 \cite{incoterms2020} 为准，
欧盟增值税申报可参考 \cite{eu_oss}。
凡涉及费率、关税、认证清单与平台仓储政策的内容，\textbf{一律以官方最新公告与费率表为准}；
本章给出的所有数字都是\kw{示范量级}，用于说明数量级与拆解结构，不可直接当作报价或成本表使用。

\section{选品方法论：把直觉变成可复核的闸门}
\label{sec:sup-select}

\subsection{为什么「找到爆款」是一个错误的提法}

新手最常问的问题是「现在什么好卖」。这个问题有三个内在缺陷。
第一，\textbf{「好卖」的信息一旦公开，超额利润就开始消失}：
平台榜单、短视频热榜、行业媒体的选品报告都是公开信息，
看到的人越多，供给进入越快，价格带下移越快。
第二，\textbf{「好卖」不等于「你能卖」}：同一个类目里，头部卖家的优势可能是评论存量、品牌认知、供应链独家款，
而你只有新品期的广告预算。
第三，也是最重要的：\textbf{选品不是一次判断，而是一串判断}。
一个商品从想法到上架要经过需求、竞争、利润、供应链、打样、复购六道关口，
每一道都会淘汰掉一半以上的候选。把精力花在「想到一个好点子」上，而不是花在「建一条能快速否决坏点子的流水线」上，
是选品效率低下的根本原因。

\kw{闸门思维}的产物是一条漏斗：假设每轮评估 40 个候选（示范量级），
需求验证淘汰 60\%，竞争分析再淘汰 50\%，利润测算再淘汰 50\%，供应链可行性再淘汰 50\%，
最终进入打样的通常只有 1--3 个。这条漏斗的意义在于：\textbf{否决的成本远低于试错的成本}。
在纸上否决一个候选的成本是半小时，在仓库里否决一个候选的成本是几万到几十万元。

\begin{figure}[htbp]
\centering
\begin{tikzpicture}[
  node distance=3mm,
  stage/.style={draw=NavyMid, fill=BoxBG, rounded corners=2pt, minimum height=10mm,
                inner sep=3pt, text width=1.72cm, align=center, font=\sffamily\scriptsize},
  gate/.style={draw=AmberDeep, fill=white, rounded corners=1pt, minimum height=10mm,
               inner sep=2pt, text width=1.72cm, align=center, font=\sffamily\tiny}
]
\node[stage] (a) at (0,0) {\textbf{需求验证}\\搜索量/趋势/差评};
\node[gate, right=of a] (g1) {闸门一\\需求是否真实};
\node[stage, right=of g1] (b) {\textbf{竞争分析}\\评论/价格带/广告};
\node[gate, right=of b] (g2) {闸门二\\能否进前 20};
\node[stage, right=of g2] (c) {\textbf{利润测算}\\全链路单位经济};
\node[gate, right=of c] (g3) {闸门三\\净利率是否达标};
\node[stage] (d) at ([yshift=-17mm]a) {\textbf{供应链可行性}\\MOQ/交期/定制};
\node[gate, right=of d] (g4) {闸门四\\起订资金可否承受};
\node[stage, right=of g4] (e) {\textbf{打样与小批量试单}\\用真实数据验证};
\node[gate, right=of e] (g5) {闸门五\\退货与复购是否可控};
\node[stage, right=of g5] (f) {\textbf{上架与复盘}\\滚动迭代与下架};
\draw[-{Stealth}, thick, NavyMid] (a) -- (g1);
\draw[-{Stealth}, thick, NavyMid] (g1) -- (b);
\draw[-{Stealth}, thick, NavyMid] (b) -- (g2);
\draw[-{Stealth}, thick, NavyMid] (g2) -- (c);
\draw[-{Stealth}, thick, NavyMid] (c) -- (g3);
\draw[-{Stealth}, thick, AmberDeep] (g3.south) -- ++(0,-5mm) -| (d.north);
\node[AmberDeep, font=\sffamily\tiny, below] at (5.9,-1.05) {通过后进入下一轮};
\draw[-{Stealth}, thick, NavyMid] (d) -- (g4);
\draw[-{Stealth}, thick, NavyMid] (g4) -- (e);
\draw[-{Stealth}, thick, NavyMid] (e) -- (g5);
\draw[-{Stealth}, thick, NavyMid] (g5) -- (f);
\node[draw=RuleGray, fill=BoxBGAlt, rounded corners=2pt, inner sep=3pt,
      font=\sffamily\tiny, text width=11.6cm, align=left] at (5.9,-3.05)
  {每一道闸门都必须有\emph{可量化}的通过标准与\emph{明确的否决线}。没有否决线的闸门等于没有闸门：
   团队会不断用「再观察一下」把坏候选推进到下一轮，最后所有成本都压在打样与备货上。};
\end{tikzpicture}
\caption{选品到上架的流程与五道闸门：每一道闸门都是一次低成本否决的机会}
\label{fig:sup-gates}
\end{figure}

\subsection{需求验证：四种信号与它们的可信度}

需求验证的核心问题是「\emph{是否真的有一群人，正在用钱表达这个需求}」。
可用的信号有四种，可信度依次递减，必须交叉验证。

\textbf{(1) 平台站内搜索量}。Amazon 的 Brand Analytics 搜索词报告给出\kw{搜索频率排名}（Search Frequency Rank），
它是相对排名而非绝对次数，需要用若干已知词做锚点才能换算成量级。
第三方工具（Helium 10、卖家精灵、Jungle Scout 等）给出的「月搜索量」是估算值，
不同工具之间常有 30\%--50\% 的差异（示范量级）。
经验判断：一个细分词的月搜索量在 \textbf{5 万--20 万}属于中腰部可做的区间；
低于 3000 且客单价低于 20 美元的词，除非能靠关联流量或复购撑起来，否则不值得单独开品。

\textbf{(2) Google Trends 与外部需求趋势}。Trends 给出的是 0--100 的\emph{相对}指数，
它的正确用法不是看绝对值，而是看三件事：
\emph{五年趋势的斜率}（长期增长还是一次性爆款）、\emph{季节性的形状与峰值月份}（决定备货节奏）、
以及\emph{相关查询里的飙升词}（发现需求正在迁移的方向）。
一个典型的失败用法是把某次短视频带火的尖峰当成基线，按峰值备货，结果三个月后需求回到原点，
库存砸在手里——这正是本章后面要讲的典型陷阱。

\textbf{(3) 平台榜单}。Amazon Best Sellers、Movers \& Shakers、New Releases，TikTok Shop 的热销榜，
以及 Temu 与 SHEIN 的类目排行 \cite{temu_about,shein_about}。
榜单的价值在于它验证了「这个需求已经被平台流量验证过」，
代价是它同时也说明了「供给已经就位」。
更实用的读法是看 \emph{New Releases} 与 \emph{Movers \& Shakers}：
前者反映新品能不能起来，后者反映短周期的需求突变。
如果某个细分里长期只有老品在榜、新品榜常年空缺，通常意味着类目已被评论存量锁死。

\textbf{(4) 评论挖掘}。这是四种信号里信息密度最高、也最需要动手的一种：
把竞品的中差评（1--3 星）与高赞好评各抓几百条，做痛点聚类。
高频出现的抱怨就是产品改进清单，高频出现的赞美就是必须保留的设计特征。
一个可操作的经验：如果前 20 条差评里有 3 条以上指向同一个可解决的结构性问题
（例如「用两周就断连」「充电盒盖子松」「尺寸偏小两码」），
那么这个类目就存在\kw{差异化切口}；反之，如果差评全是「物流慢」「包装破损」这类与产品无关的抱怨，
说明产品本身已经同质化，切进去只能打价格战。

\begin{keypoint}[需求验证的输出物是一张带否决线的判断表，不是一句结论]
一个成熟团队的需求验证交付物通常包含四列：\emph{信号来源}、\emph{观测值}、\emph{判定阈值}、\emph{结论}。
例如：站内月搜索量 8.6 万（阈值 5 万以上）、Trends 五年斜率向上且在 11--12 月出现峰值（用于排备货节奏）、
Top 20 中近 12 个月新上架且进入前 20 的链接有 3 条（说明评论存量未锁死）、
差评聚类出 2 个可解决的结构性问题（说明存在差异化切口）。
四条同时命中才放行。\textbf{任何一条不达标就否决，而不是「综合考量后仍然推进」}——
闸门的作用就是在信息还不完整时强制做出否决，因为继续投入只会让否决变得更贵。
\end{keypoint}

\subsection{竞争分析：评论存量、价格带与广告密度}

需求真实不等于你打得进去。竞争分析要回答三个可量化的问题：
\emph{进入壁垒有多高}、\emph{价格空间在哪里}、\emph{获客有多贵}。

\textbf{评论存量与集中度}。把类目 Top 20 的评论数按降序排列，观察三件事：
第 1 名的评论数（头部壁垒）、第 10 名的评论数（进入门槛）、
以及前 5 名占 Top 20 总评论数的比例（集中度 $CR_5$）。
经验阈值（示范量级）：如果第 1 名评论数超过 2 万且客单价低于 25 美元，正面进攻基本不可行；
如果第 10 名评论数在 500--2000 之间，说明类目仍在洗牌期，新品有机会；
如果 $CR_5 > 0.7$，说明流量被少数链接垄断，新进入者需要靠站外流量或差异化款式绕开正面战场。

\textbf{价格带分布}。把 Top 50 的售价做成直方图，找「\kw{价格空档}」：
某个区间的链接数量明显低于两侧，且这个区间能覆盖成本。
价格空档有两种成因：一种是成本结构导致的（这个价位做不出合格的品质），
另一种是认知惯性导致的（大家都认为应该卖 19.99 或 39.99）。
前者是陷阱，后者是机会——但要用小批量试单去验证消费者是否接受这个价位。

\textbf{广告密度}。搜索结果首页（前 20 位）中 Sponsored 位置的数量是最直接的获客成本信号。
如果首页有 12 个以上广告位被买断，说明自然位的曝光被严重挤压，
新品必须靠广告起步，且单次点击成本会被抬得很高。
配合第三方工具的估算 CPC 与预计 $\CVR$，可以先算出「\emph{广告盈亏平衡的转化率}」：
在给定售价与毛利率下，允许的最高单次获客成本决定了你能否承受这个类目的竞价水平。
这部分逻辑与 \cref{ch:ad,ch:meta} 里的投放框架是同一条链，只是站在品类选择的角度提前做了一次估算。

\subsection{利润测算：全链路单位经济}

这是选品环节最容易出错、也最容易补救的一步。绝大多数「看起来赚钱」的失败，
根因都是只算了\emph{采购价与售价的差额}，而漏掉了下面这七项中的某几项。

\begin{equation}
\begin{aligned}
\pi_{\text{单位}} ={}& P_{\text{售价}} - C_{\text{出厂}} - F_{\text{头程}} - D_{\text{关税}} - C_{\text{平台佣金}}
 - F_{\text{履约}} \\
 &- A_{\text{广告}} - r_{\text{退货}}\cdot L_{\text{退货}} - O_{\text{其他}} ,
\end{aligned}
\label{eq:unitecon}
\end{equation}

其中 $r_{\text{退货}}$ 是退货率，$L_{\text{退货}}$ 是\emph{单笔退货的净损失}
（逆向运费 + 检测翻新 + 残值折损，见 \cref{eq:return}），
$O_{\text{其他}}$ 包含仓储费、长期仓储或库龄附加费、支付手续费、汇率损耗与合规摊销。
\emph{注意退货项是乘性的}：退货率从 5\% 涨到 25\%，即使其他项完全不变，
毛利也可能被吃掉一半以上。这解释了为什么服装类目的选品标准与 3C 完全不同。

\begin{table}[htbp]
\centering
\caption{一款售价 29.99 美元的小件电子产品的全链路单位经济（示范量级，非真实报价）}
\label{tab:sup-profit}
\small
\begin{tabularx}{\textwidth}{L{2.9cm} R{1.9cm} R{1.6cm} Y}
\toprule
\tabhead{成本项} & \tabhead{金额（美元）} & \tabhead{占售价} & \tabhead{口径与常见误读} \\
\midrule
售价 $P$ & 29.99 & 100\% & 含税价；实际到手还要扣平台佣金与履约 \\
出厂价 $C$ & 8.60 & 28.7\% & FOB 中国港口；不含包装唛头与质检费 \\
头程 $F$ & 1.35 & 4.5\% & 海运散货摊到单件的费用，按体积重或 CBM 计价 \\
关税 $D$ & 2.25 & 7.5\% & 按 HS 编码税率乘申报价值；以官方最新税则为准 \\
平台佣金 & 4.50 & 15.0\% & 多数类目在 8\%--20\% 区间，另有支付与广告附加 \\
履约费 & 4.10 & 13.7\% & 平台仓或第三方仓的拣配打包派送，按尺寸分段 \\
广告 $A$ & 4.20 & 14.0\% & 成熟期总广告成本占收入比的经验区间为 10\%--20\% \\
退货损失 & 0.90 & 3.0\% & 退货率 6\% 乘单笔净损失约 15 美元 \\
其他 $O$ & 1.29 & 4.3\% & 仓储、长期仓储、支付、汇率与合规摊销 \\
\midrule
\textbf{单位净利} & \textbf{2.80} & \textbf{9.3\%} & 低于 5\% 时，一次汇率波动或一次关税调整即可转负 \\
\bottomrule
\end{tabularx}
\end{table}

从 \cref{tab:sup-profit} 可以读出一个反直觉的结论：\textbf{出厂价只占售价的不到三成}，
而\emph{履约 + 广告 + 退货}三项合计接近售价的三分之一。
这意味着选品时真正要盯的不是「能不能把出厂价从 8.6 压到 8.0」（省 2 个百分点），
而是「这个商品的体积重是不是过大」（一次就能吃掉 5--10 个百分点）。
这也是本章陷阱盒要展开的主题。

\subsection{选品打分模型：权重从哪来}

多因子加权打分是选品的标准工具，但大部分打分表死在同一个地方：\textbf{权重是拍的}。
一旦权重靠拍，打分表就退化成「把已经决定要做的品包装成数据支持的样子」。
正确的权重来源有三条路径，按可信度排序：

\textbf{(1) 用历史数据回归。} 把过去 12--24 个月已上架 SKU 的每个因子值作为特征，
用「90 天售罄率」「毛利率达成」「180 天动销率」作为标签，做逻辑回归或梯度提升，
得到的系数即权重。这条路径的前提是 SKU 数量足够（经验上 100 个以上才有意义），
且因子在历史上有过足够的变化幅度。

\textbf{(2) 用专家判断做结构化聚合。} 让 5--7 位不同岗位（采购、运营、投放、财务、客服）的人独立打分，
取几何平均而不是算术平均——几何平均会惩罚「某一维度极差」的候选，这正是我们想要的。
再用一致性检验剔除明显离群的判断。

\textbf{(3) 用敏感性分析代替精确权重。} 如果既没有数据也没有足够专家，
就把权重设成区间，然后检验「排序是否稳定」：
如果候选 A 在权重扰动 $\pm 20\%$ 的范围内始终排第一，就可以放行；
如果排序频繁翻转，说明这个决策的信息量根本不足以支撑判断，应该退回需求验证阶段。

打分模型的形式化定义如下。设 $m$ 个因子，第 $i$ 个因子的权重为 $w_i$，
候选商品 $j$ 在第 $i$ 个因子上的原始取值为 $s_{ij}$：

\begin{equation}
\tilde{s}_{ij} = \frac{s_{ij} - \min_k s_{kj}}{\max_k s_{kj} - \min_k s_{kj}} \in [0,1],
\qquad
\text{Score}_j = 100\sum_{i=1}^{m} w_i\,\tilde{s}_{ij},
\qquad
\sum_{i=1}^{m} w_i = 1 .
\label{eq:scorecard}
\end{equation}

反向因子（越小越好，例如竞争强度、体积重、退货率、交期）取 $\tilde{s}^{\,\prime}_{ij} = 1 - \tilde{s}_{ij}$。
几个工程细节值得强调：
\emph{归一化必须在候选池内做}，而不是用全类目的极值，否则同一批候选的分数不可比；
\emph{权重之和必须为 1}，否则总分不可解释；
\emph{打分表要留一列「数据来源」}，凡是拿不到来源的因子必须降权或删除，
否则整张表会退化成主观印象的加权平均。

\begin{table}[htbp]
\centering
\caption{选品加权打分表：因子、权重区间与判定阈值（权重为示范量级，须用自有数据回归校准）}
\label{tab:sup-score}
\small
\begin{tabularx}{\textwidth}{L{2.3cm} C{1.2cm} R{1.3cm} L{2.6cm} Y}
\toprule
\tabhead{因子} & \tabhead{方向} & \tabhead{权重} & \tabhead{数据来源} & \tabhead{判定阈值（示范量级）} \\
\midrule
需求强度 & 正向 & 0.22 & 站内搜索量、Trends、平台榜单 & 月搜索量 $>$ 5 万，五年趋势不向下 \\
竞争强度 & 反向 & 0.18 & Top 20 评论数、$CR_5$、广告位占比 & 第 10 名评论数 $<$ 2000 \\
毛利空间 & 正向 & 0.20 & \cref{eq:unitecon} 全链路测算 & 净利率 $>$ 10\%，毛利率 $>$ 45\% \\
物流友好度 & 正向 & 0.12 & 体积重、易碎性、带电性 & 体积重与实重之比 $<$ 2，无空运限制 \\
合规门槛 & 正向 & 0.08 & 认证清单与测试周期 & 无强制认证，或认证周期 $<$ 6 周 \\
现金流周期 & 反向 & 0.10 & 生产 + 在途 + 结算账期 & 现金周期 $<$ 120 天 \\
供应链可控性 & 正向 & 0.10 & 供应商数量、可定制性 & 存在 2 家以上可比供应商 \\
\midrule
\multicolumn{5}{@{}>{\raggedright\arraybackslash}p{\linewidth}@{}}{\textbf{决策规则：总分 $\ge$ 70 进入打样；55--70 放入观察池；低于 55 直接否决}}\\
\bottomrule
\end{tabularx}
\end{table}

\subsection{供应链可行性：MOQ、起订资金、交期与定制能力}

打分表上「毛利空间」这一项，最终能不能兑现，取决于供应链是否配合。
四个必须提前问清的量：

\textbf{(1) MOQ（最小起订量）}。同一个款式的不同颜色、不同尺码通常各自有 MOQ。
一个常见的坑是：业务希望上架 6 个颜色，供应商报的 MOQ 是「每色 500 件」，
于是首单变成了 3000 件而非 500 件。必须把\emph{总 MOQ} 与\emph{单 SKU MOQ} 分开确认。

\textbf{(2) 起订资金与现金周期}。首单资金占用不只是货款：

\begin{equation}
K_{\text{首单}} = \text{MOQ}\cdot C_{\text{出厂}} + M_{\text{模具开版}} + T_{\text{打样}} + Q_{\text{检测认证}},
\label{eq:moqcash}
\end{equation}

\begin{equation}
T_{\text{现金}} = T_{\text{生产}} + T_{\text{在途}} + T_{\text{入仓上架}} + T_{\text{平台结算}} - T_{\text{供应商账期}} .
\label{eq:cashcycle}
\end{equation}

按示范量级代入：生产 30 天、海运在途 32 天、入仓上架 10 天、平台结算 14 天、供应商账期 0 天，
则现金周期 $T_{\text{现金}} \approx 86$ 天。如果按 4 个月备货节奏滚动，需要的营运资金约为单批货值的 2--3 倍。
\emph{很多「毛利率很好」的选品，最终死在这里}：不是不赚钱，而是钱都被在途和库存占住了，没有钱投下一批。

\textbf{(3) 交期与交期波动}。要区分三个数字：\emph{常规交期}（正常排产下的天数）、
\emph{旺季交期}（节前 2--3 个月，通常延长 50\%--100\%）、
以及\emph{实际交期达成率}（历史订单中按期出货的比例）。
只有第三个数字才有资格进入安全库存公式里的 $\sigma_{LT}$。供应商口头承诺的「30 天」几乎从不包含旺季。

\textbf{(4) 定制能力}。三个层次：\kw{白牌}（贴你的标，无设计改动）、
\kw{OEM}（按你的图纸生产，需要开模费与更长的打样周期）、
\kw{ODM}（供应商提供设计方案，你选款，独家性最弱）。
定制程度越高，独家性越强、溢价空间越大，但起订资金、开发周期与失败成本也同步上升。
对新品类，比较稳的路径是先在白牌上验证需求（与 \cite{ries2011lean} 的最小可行产品思路一致），
确认有稳定复购后再投入开模做独家款。

\subsection{案例：一款蓝牙耳机的六周评估与半年上市}

\begin{casebox}[案例：一款真无线耳机从选品到上市的成本与周期（示范量级）]
这是一个依据行业公开通识构造的示范复盘，不代表任何具体公司的内部数据；所有金额与周期均为量级示意。

\textbf{背景}：某卖家计划上架一款真无线蓝牙耳机（TWS），目标市场美国，售价定位 29.99 美元。

\textbf{评估阶段（约 6 周）}：
\begin{enumerate}
  \item 第 1--2 周，需求验证。站内主词月搜索量约 40 万，长尾词「open ear clip」出现明显上升趋势，
        说明需求正在从入耳式向开放式迁移；抓取 400 条中差评，聚类出三个高频抱怨：
        久戴疼痛、通话降噪差、充电盒易松动。
  \item 第 3 周，竞争分析。Top 20 中第 1 名评论数 8 万，但第 10 名仅 1200，$CR_5 \approx 0.55$，
        首页广告位 11 个，估算单次点击成本约 1.1 美元。结论：不正面打头部爆款，切开放式细分。
  \item 第 4 周，利润测算。按 \cref{eq:unitecon} 逐项代入，得到净利率约 9\%（\cref{tab:sup-profit}）。
        敏感性检查发现最大风险不是出厂价，而是体积重与旺季广告成本。
  \item 第 5--6 周，供应链可行性。联系 7 家供应商，2 家为工厂、5 家为贸易商；
        工厂报价 MOQ 每色 1000 件、交期 30 天（旺季 55 天）；起订资金按 \cref{eq:moqcash} 测算约 12 万元人民币，
        现金周期约 88 天。
\end{enumerate}

\textbf{执行阶段（约 5 个月）}：打样 3 轮（外观样、功能样、量产样，共 4 周，打样费 3000 元量产后返还）
$\to$ 第三方验厂（约 6000 元，含质量体系与产线核对）$\to$ 首单 2000 件（生产 32 天）
$\to$ 海运散货（在途 33 天，摊到单件约 1.35 美元）$\to$ 入仓上架（约 10 天）。

\textbf{结果形态}：上架后第 1 个月靠广告起量，广告成本占收入比一度达到 60\% 以上（亏损期，属正常）；
第 3 个月随着评论累积到 200 条、自然位进入前 20，该比例降到 13\% 左右，单件净利回到 2.5--3.0 美元区间。
第 6 个月复盘时发现真正的利润来源不是耳机本体，而是\emph{耳挂配件与备用耳塞的组合销售}——
它的体积重极低、退货率不到 2\%，是典型的「高毛利补货型 SKU」。

\textbf{可迁移的结论}：选品的价值不在选中了哪一个品，而在\emph{用六周的低成本评估，换掉六个月的不可逆投入}。
同时，真正赚钱的往往是主品带出来的配件与耗材——评估时就要把「复购与关联销售」当成一个独立因子来打分。
\end{casebox}

\begin{pitfall}[只按「毛利率高」选品，被体积重与关税吃掉全部利润]
\textbf{现象}：一款家居抱枕的采购价 12 元人民币，售价 29.99 美元，纸面毛利率超过 80\%，
团队判定为「超高毛利品」并首批备货 5000 件。上架三个月后财务报表显示单件净亏约 1.2 美元。

\textbf{根因}（四个叠加项，任何一项单独出现都不致命）：
\begin{enumerate}
  \item \textbf{体积重被忽略}。抱枕实重 0.4 kg，但包装后尺寸约 $45\times35\times18$ cm，
        按体积重系数 6000 计算得 4.7 kg，是实重的近 12 倍。
        头程与履约费全部按 4.7 kg 计费，单件履约成本从预估的 3 美元涨到 9 美元以上。
  \item \textbf{关税按品类与申报价值计，而不是按采购成本计}。
        关税基数是申报价值（成交价），与「采购价低」无关；
        低采购价并不带来低关税，只带来更小的毛利基数来吸收同一笔关税。
  \item \textbf{退货率被低估}。家居软装类目的退货率常在 5\%--15\%（示范量级），
        而抱枕退货后无法二次销售（卫生与压痕问题），残值回收率接近 0，
        使 \cref{eq:return} 里的净损失接近全额货值。
  \item \textbf{长期仓储费按体积计}。体积大、周转慢的商品在旺季后的仓储附加费按体积计，
        滞销三个月后，仓储费已经接近甚至超过货值本身。
\end{enumerate}

\textbf{排查方法}：
(1) 在利润测算表里\emph{强制加入体积重字段}，并计算「体积重与实重之比」，比值超过 2 的候选自动降权；
(2) 用「售价除以体积重」作为关键比率，行业经验上低于 6 美元/kg 的候选要格外谨慎；
(3) 关税一律按 HS 编码查官方最新税则，禁止用「大概 5\%」估算；
(4) 对退货后不可二次销售的品类，把退货率与残值回收率同时显式建模，
而不是用一个笼统的「退货率 5\%」；
(5) 上线后设立「实际到账净利与测算净利」的月度对账，偏差超过 20\% 就必须回头修正测算模板。
\end{pitfall}

\section{供应商管理：从找源到合同}
\label{sec:sup-supplier}

\subsection{找源：1688、义乌与产业带}

供应商来源按「距离工厂的远近」分成三层，成本与风险依次变化。

\textbf{第一层：线上平台（1688、阿里国际站、Made-in-China）。}
1688 的优势是选择多、比价快，劣势是\emph{贸易商与工厂混在一起}。
快速鉴别工厂的三个问题：
「你们有几条产线、日产能多少」、
「能不能现在拍一段车间视频给我看」（贸易商通常拖延或转发别人的素材）、
「打样费和大货单价的差额是多少」（工厂的打样费往往更接近成本，且量产后能返还）。
平台上的「实力商家」「深度验商」「回头率」「近 90 天成交额」是可以直接看的信号，
但要注意：\emph{回头率高}说明老客户续单多，通常意味着配合度好；
\emph{成交额高但回头率低}可能是靠一次性低价冲量的贸易商。

\textbf{第二层：专业市场（义乌国际商贸城、广州十三行与白马、深圳华强北）。}
专业市场的价值在于\kw{款式密度}与\kw{起订弹性}：
在义乌可以用极低的起订量（几十件）采购到大量小商品，
非常适合用来做「款式测试」——先用小单验证哪一个款式能打，再回头找工厂做量产。
代价是单位成本偏高、品质一致性差、且供应商极度同质化，容易被竞争对手用同一款打价格战。

\textbf{第三层：产业带（源头工厂集群）。}
这是成本最低、配合度最高、也是壁垒最强的一层。国内代表性产业带（示意）：
深圳与东莞的 3C 与配件、广州与佛山的服饰与家具、
义乌与永康的小商品与五金、南通的蓝印花与家纺、汕头的玩具、
泉州的运动鞋与服装、温州的眼镜与打火机、中山的灯具、宁波的小家电。
进入产业带的方法有三条：
参加广交会、华交会、CES 等展会做第一轮接触；
通过已合作的供应商介绍上下游（「同乡带同乡」在产业带内极为普遍）；
以及直接驻地拜访。产业带拜访的价值不只是价格，
更重要的是你能在一周内看到\emph{同一款产品在十家工厂的不同做法}，
这对理解成本结构极其有帮助。

\begin{toolbox}[供应商评估与验厂清单（可直接当作打分表使用）]
\textbf{A. 资质与背景（每项 0--5 分）}
\begin{itemize}
  \item 营业执照与经营范围是否覆盖该品类；是否有自营进出口资质；
  \item 是否为直接生产厂家（要求提供车间视频、设备清单、产能说明）；
  \item 已合作的品牌客户类型（能否提供合作证明或授权书）；
  \item 是否具备目标市场的认证经验（CPC、CE、FCC、FDA 等）。
\end{itemize}
\textbf{B. 产能与柔性}
\begin{itemize}
  \item 常规月产能与峰值月产能；旺季（节前 2--3 个月）排产是否会被大客户挤占；
  \item 最小起订量（区分总 MOQ 与单 SKU MOQ）、能否接受首单减半试产；
  \item 换线时间与多颜色并行生产能力。
\end{itemize}
\textbf{C. 品控体系}
\begin{itemize}
  \item 是否有独立 QC 团队与 IQC/IPQC/OQC 三段的原始记录表（要求看原始记录，不看模板）；
  \item 是否有检测设备（耐压、老化、跌落、盐雾等），还是纯靠目检；
  \item 近 6 个月的实际不良率与客诉率，以及对应的改善记录。
\end{itemize}
\textbf{D. 交期与配合度}
\begin{itemize}
  \item 近 12 个月订单的按期出货率（不是口头承诺）；
  \item 响应速度（工作日 4 小时内是否回复、能否提供中英文沟通）；
  \item 是否愿意配合打样、改包装、贴标、抽检与第三方验货。
\end{itemize}
\textbf{E. 商务条件}
\begin{itemize}
  \item 阶梯报价表、付款方式（电汇 30\% 定金较常见）、账期（30/60/90 天）；
  \item 是否承担质量赔付、返工费用与退运费用；
  \item 价格有效期与调价机制（原材料波动如何传导）。
\end{itemize}
\textbf{评分规则}：五个维度各取平均分，总分低于 3.5 分不进入打样；
\textbf{同一关键词至少保留 2 家达标的供应商}。
验厂方面，第三方审核（SGS、BV、TUV、Intertek 等）单次费用通常在几千到一万多元人民币量级，
社会责任审核（BSCI、Sedex）与质量体系审核（ISO 9001）是两种不同的东西，
报价前必须说清审什么；具体费用与排期以机构的官方报价为准。
\end{toolbox}

\subsection{打样、验厂与首单}

\kw{打样}通常要跑三轮，每一轮的目的不同，不能合并：
\emph{外观样}确认造型、颜色、表面处理与包装视觉效果；
\emph{功能样}确认电气性能、结构与装配公差；
\emph{量产样}（也叫金样或生产样）从量产线上随机抽取，
确认「批量生产的品质与打样一致」——这一步最容易被跳过，也最容易出问题。
打样费通常几百到几千元（示范量级），量产后部分供应商会返还，
谈条件时应把「量产样免费」写进报价单。

\kw{验厂}分自查与第三方两种。创业早期建议先用自查清单（见上面的工具箱）跑一遍，
在订单金额超过一定阈值或进入受监管品类（儿童用品、食品接触、医疗器械）时再上第三方。
验厂的真正价值不在拿到一张证书，而在于\emph{让供应商知道你会查}——
被查过一次的供应商在后续订单上的配合度通常明显更高。

\subsection{阶梯报价、谈判与多供应商备份}

阶梯报价的本质是供应商把\emph{固定成本与学习曲线}摊到更大的批量上，
其形状通常可以用一条幂律近似：

\begin{equation}
C(q) = C_0\left(\frac{q}{q_0}\right)^{-\alpha}, \qquad \alpha \in [0.05, 0.15] \ (\text{经验值}),
\label{eq:tierprice}
\end{equation}

其中 $q_0$ 是基准批量，$\alpha$ 是折扣弹性。$\alpha = 0.1$ 意味着批量翻倍时单价下降约 6.7\%。
这个式子的实用价值有两个：
一是\emph{判断报价是否合理}（如果供应商给出的阶梯折扣远陡于 $\alpha = 0.15$，
说明基准价虚高，或者大货规格与打样规格不一致）；
二是\emph{估算谈判空间}（批量从 1000 提到 3000，理论降价空间约 11\%，
这就是你手里的筹码量级）。

谈判的关键认知是：\textbf{价格只是众多可交换变量中的一个}。
可交换的变量至少包括：账期、MOQ、交期、包装与唛头、质检标准与 AQL 水平、
付款方式、独家款式保护期、模具归属、退换与赔付条款。
一个成熟的谈判策略是\emph{用不那么值钱的变量换值钱的变量}：
例如接受 60 天账期（对现金流价值极高）而放弃 3\% 的价格折让，
或者承诺全年最低采购量以换取旺季优先排产权。

\kw{多供应商备份}的工程做法是「主备比加活单」：
主供应商承担 70\% 的订单，备份供应商承担 30\%；
备份供应商必须\emph{持续有活单}（哪怕每月只有很小的量），
否则一旦主供应商出问题，备份供应商的模具、原料、排产都要从零开始，救不了急。
备份供应商的认证、打样与 AQL 标准应与主供应商完全一致，
并在主供应商的年度评审中把「备份启用演练」当成一项常规动作。

\begin{table}[htbp]
\centering
\caption{采购合同的关键条款清单（跨境采购场景，具体条款须经法律审核）}
\label{tab:sup-contract}
\small
\begin{tabularx}{\textwidth}{L{2.4cm} L{3.4cm} Y}
\toprule
\tabhead{条款} & \tabhead{应写清的内容} & \tabhead{常见缺失与后果} \\
\midrule
规格与签样 & 规格书作为合同附件；封存金样并由双方签认 & 只写「按样品」，出争议时无对标物，无法判定谁违约 \\
质量标准 & 引用 AQL 标准与检验水平，明确致命缺陷零容忍 & 只写「质量合格」，等于没有标准 \\
交期与违约 & 明确交货日、允许延期天数、延期违约金比例与上限 & 未约定违约金或约定过低，延期无成本 \\
价格与调价 & 阶梯价表、价格有效期、原材料波动的传导机制 & 无有效期，旺季被临时加价 \\
付款方式 & 定金比例、尾款条件、验货通过后付款 & 全额预付，丧失全部质量杠杆 \\
赔付与退换 & 不良品的补货、返工、退运费用承担方 & 只写「协商解决」，实际承担方是你 \\
模具与知识产权 & 模具归属、独家款式保护期、图纸保密义务 & 未约定模具归属，换供应商时被卡住 \\
贸易术语 & 明确 EXW/FOB/CIF/DDP 等术语及适用版本 & 术语不清，运费与风险划分扯皮 \\
争议解决 & 适用法律与仲裁机构（如 CIETAC） & 未约定，跨境诉讼成本极高 \\
\bottomrule
\end{tabularx}
\end{table}

\cref{tab:sup-contract} 里的\kw{贸易术语}值得单独说明。
按 ICC 的 \cite{incoterms2020}，\code{EXW} 是工厂交货，买方承担全部运输与出口手续；
\code{FOB} 是装运港船上交货，卖方负责出口清关与装船前的费用，这是跨境采购最常用的术语；
\code{CIF} 由卖方承担运费与保险到目的港；
\code{DDP} 由卖方承担目的港关税与派送，对买方最省事但报价最高。
术语的选择直接决定了关税与头程由谁申报、由谁承担，\emph{在谈价之前就必须定下来}，
否则同一批货在两个术语下的到手成本可能相差 15\% 以上。

\section{品控与合规}
\label{sec:sup-qc}

\subsection{三段检验：IQC、IPQC、OQC 与 AQL 判定}

品控不是「最后验一次货」，而是分布在三个位置的三道关卡：

\textbf{IQC（来料检验）}：原材料与零部件入库前的检验。
跨境的特殊性在于，很多组件（电芯、芯片、面料）来自你的上游供应商，
一旦它们的批次有问题，成品的所有检验都是白做。
IQC 的重点是\emph{批次一致性}与\emph{关键参数}，而不是全检所有指标。

\textbf{IPQC（过程检验）}：生产过程中的巡检与首件确认。
核心动作是\emph{首件检验}（每个批次开机后的第一件必须全项检验并留样）
与\emph{按固定频率抽检}。IPQC 的价值在于把问题在\emph{增值之前}拦住：
一件在注塑工序就被发现缩水的壳体，成本是几毛钱；
装成整机再发现，成本是几十元。

\textbf{OQC（出货检验）}：成品入库或装柜前的最终检验，通常按 AQL 抽样标准执行。
国际贸易中最常引用的标准是 ISO 2859-1（国标 GB/T 2828.1），
采用\emph{一般检验水平 II、正常检验、一次抽样}时，
「批量大小」决定「样本量字码」，字码决定「样本量 $n$」，
再由 AQL 值查出接收数 $Ac$ 与拒收数 $Re$。

抽样的统计含义是：\textbf{给定一个真实不良率 $p$，这批货被接收的概率}可以用超几何分布精确计算，
工程上用二项分布近似：

\begin{equation}
\begin{aligned}
L(p) ={}& \Prob(\text{接收}\mid p) = \sum_{d=0}^{Ac} \binom{n}{d} p^{d}(1-p)^{n-d}, \\
&\text{判定：}\ d \le Ac \ \text{接收};\quad d \ge Re \ \text{拒收},\quad Re = Ac + 1 .
\end{aligned}
\label{eq:aql}
\end{equation}

这条公式能回答一个极其常见但很少有人算的问题：
\emph{「AQL 2.5」到底放过了多少不良品？}
在 $n = 200$、$Ac = 10$ 的设置下（示范量级，对应较大的批量），
若真实不良率 $p = 2\%$，接收概率约为 0.99；
若 $p = 5\%$，接收概率仍有约 0.58；若 $p = 8\%$，接收概率降至约 0.15。
换句话说：\textbf{AQL 不是「不良率上限」，而是「验收规则」}。
一批真实不良率 5\% 的货，仍有一半以上的概率被放行。
业务上必须接受这个统计事实，并用\emph{额外条款}（如「客户投诉率超过 X\% 由供应商赔付」）
来补足抽样检验拦不住的长尾风险。

\begin{table}[htbp]
\centering
\caption{一次抽样正常检验的对照示例（一般检验水平 II，简化示意，正式使用须查标准原表）}
\label{tab:sup-aql}
\small
\begin{tabularx}{\textwidth}{R{2.0cm} C{1.8cm} C{1.5cm} C{2.0cm} C{2.0cm} Y}
\toprule
\tabhead{批量范围} & \tabhead{样本量字码} & \tabhead{样本量 $n$} & \tabhead{主要缺陷 AQL 2.5} & \tabhead{次要缺陷 AQL 4.0} & \tabhead{备注} \\
\midrule
501--1200 & J & 80 & $Ac=5,\ Re=6$ & $Ac=7,\ Re=8$ & 中小批量首单常见区间 \\
1201--3200 & K & 125 & $Ac=7,\ Re=8$ & $Ac=10,\ Re=11$ & 首单主力区间 \\
3201--10000 & L & 200 & $Ac=10,\ Re=11$ & $Ac=14,\ Re=15$ & 返单常见区间 \\
10001--35000 & M & 315 & $Ac=14,\ Re=15$ & $Ac=21,\ Re=22$ & 大批量，需提前排验货资源 \\
\midrule
\multicolumn{6}{@{}>{\raggedright\arraybackslash}p{\linewidth}@{}}{\textbf{致命缺陷（安全、法规、功能失效）通常采用零容忍：$Ac=0$，发现即整批拒收}}\\
\bottomrule
\end{tabularx}
\end{table}

\begin{pitfall}[把 AQL 当成「合格率承诺」，验货流于形式]
\textbf{现象}：供应商在合同里写「按 AQL 2.5 收货」，验货时第三方出具「通过」报告，
但上架后 30 天内的退货率仍高达 12\%，差评集中在同一个功能缺陷上。

\textbf{根因}：
(1) \textbf{抽样只覆盖了「当时能看到」的缺陷}。功能失效、电池衰减、涂层脱落这类缺陷
在出厂时的暴露率很低，需要老化测试或加速试验才能暴露，标准抽样根本不覆盖；
(2) \textbf{缺陷分类被做手脚}。把功能性缺陷归类为「次要缺陷」，用 AQL 4.0 甚至 6.5 去验收，
接收数一放宽，通过率自然高；
(3) \textbf{抽样的随机性被破坏}。验货员按供应商安排的「展示区」抽样，而不是从已封箱的成品中随机抽取；
(4) \textbf{没有把市场端数据回流}。验货报告与退货原因分析是两套系统，缺陷模式无法闭环。

\textbf{排查方法}：
(1) 在合同中\emph{明确列出致命、主要、次要缺陷的定义清单}，功能性缺陷一律划入主要或致命；
(2) 对电气与耐用性要求高的品类，增加\emph{批次可靠性抽测}（如高低温循环、充放电循环），
用少量样本换取对长尾失效的早期预警；
(3) 规定抽样必须从\emph{已封箱的待发成品}中随机进行，验货员有权自行选择箱号；
(4) 建立「退货原因代码」与「缺陷类别」的映射表，每月计算「出厂 AQL 通过批」的\emph{市场退货率}，
当该指标超过阈值时，说明验收标准需要收紧，而不是供应商突然变差。
\end{pitfall}

\subsection{质量成本的四个分类}

质量管理里最实用的一张分类表，是把所有与质量相关的支出分成四类：

\begin{equation}
C_{\text{质量}} = C_{\text{预防}} + C_{\text{鉴定}} + C_{\text{内部失败}} + C_{\text{外部失败}} .
\label{eq:copq}
\end{equation}

\begin{table}[htbp]
\centering
\caption{质量成本的四个分类、典型科目与对利润的杠杆}
\label{tab:sup-copq}
\small
\begin{tabularx}{\textwidth}{L{2.3cm} L{3.9cm} L{2.4cm} Y}
\toprule
\tabhead{分类} & \tabhead{典型科目} & \tabhead{发生位置} & \tabhead{杠杆特性} \\
\midrule
预防成本 & 供应商审核、打样、标准作业流程与培训、设计与工艺验证、认证准备 & 生产之前 & 每投入 1 元，可省去后续数倍的失败成本（经验比例约 1:10:100） \\
鉴定成本 & IQC/IPQC/OQC 人工、第三方验货、实验室测试、老化测试 & 生产之中 & 随检验强度线性上升，存在「检到无利可图」的边界 \\
内部失败 & 报废、返工、重工工时、批次延期导致的空运补货 & 交付之前 & 直接吃掉毛利，但客户无感，属于「可控的坏消息」 \\
外部失败 & 客诉赔付、退货与逆向物流、差评导致的转化率下降、平台处罚、召回 & 交付之后 & 成本最高且带有\emph{声誉复利}效应，必须优先避免 \\
\bottomrule
\end{tabularx}
\end{table}

\cref{tab:sup-copq} 的关键洞察是「1:10:100」这个经验比例：
在预防阶段花 1 元解决的问题，到了内部失败阶段要花约 10 元，
到了外部失败阶段要花约 100 元（示范量级，来自质量管理通识，用于说明数量级而非精确预测）。
跨境场景会把这个比例放得更大，因为外部失败还额外包含：
逆向国际物流、无法二次销售的残值损失、以及差评对后续转化率的持续侵蚀。
经验上，商品评分从 4.5 掉到 4.0 会让同类目转化率下降 10\%--30\%（示范量级），
而这个损失是\emph{永久性}的，不会随补货而恢复。

因此，跨境品控的资源配置原则是：\textbf{把预算前移到预防与鉴定，把外部失败当作不可接受事件来管理}。
具体到操作层面：在设计阶段做可靠性验证、在下单前做供应商审核与打样、
在量产中做过程巡检、在出货前做 AQL 抽检加可靠性抽测、
在上架后建立退货原因回流机制。

\subsection{认证与合规准备}

认证是跨境选品里最容易被低估的一项：它不产生任何收入，
但\emph{没有它就没法上架}，而且周期常常长达数周到数月，直接决定上市时间。
下面的清单只作为\emph{方向性提示}，
\textbf{具体的适用法规、测试项目、发证机构与费用一律以官方最新公告与认可的检测机构报价为准}。

\begin{table}[htbp]
\centering
\caption{主要目标市场的常见认证与合规要求（方向性提示，非法律意见）}
\label{tab:sup-cert}
\small
\begin{tabularx}{\textwidth}{L{1.9cm} L{2.9cm} L{4.0cm} Y}
\toprule
\tabhead{市场} & \tabhead{常见要求} & \tabhead{适用情形} & \tabhead{周期与准备要点} \\
\midrule
美国 & CPC（儿童产品证书） & 面向 12 岁以下儿童的产品 & 须由 CPSC 认可的第三方实验室测试，周期常以周计，需留存测试报告备查 \\
美国 & FCC 认证 & 含无线发射或数字电路的电子产品 & 视产品分认证或供应商符合性声明，射频设计改动需重新测试 \\
美国 & FDA 相关注册 & 食品接触材料、化妆品、医疗器械、激光产品 & 不同类别的注册与列名要求差异极大，须按具体品类核对官方指引 \\
欧盟 & CE 标志 & 低电压、电磁兼容、无线设备、玩具、机械等指令覆盖的产品 & 需完成符合性评定、技术文件与符合性声明；部分品类须公告机构介入 \\
欧盟 & RoHS 与 REACH & 电子电气产品、含化学品的产品 & 限制物质检测，供应链需提供材质声明 \\
欧盟 & EPR 与 WEEE 及包装法 & 在成员国销售的商品 & 需在各成员国注册并申报，跨境电商可通过一站式申报处理部分增值税 \cite{eu_oss} \\
英国 & UKCA & 原 CE 覆盖的多数品类 & 与 CE 并行存在，须按英国官方最新指引确认接受规则 \\
\bottomrule
\end{tabularx}
\end{table}

关于认证，有三条来自实践的经验判断值得记住。
第一，\textbf{认证要在选品阶段就确认，而不是在下单之后}：
一个需要 FCC 认证的无线产品，射频方案的任何改动都会导致重新测试，
这会把上市时间推迟一个月以上。
第二，\textbf{要区分「强制认证」与「平台要求」}：
CE 标志是法规要求，而平台可能额外要求 UL 或 ETL 报告、特定检测报告，
两者都不能省，但预算和周期要分开规划。
第三，\textbf{认证文件是资产也是负债}：
它既是上架的通行证，也是一份随时可能被平台或监管机构抽查的法律文件，
测试报告的真实性与一致性必须严格控制，绝不能接受供应商提供的「来源不明的证书」。

\section{库存与计划}
\label{sec:sup-inventory}

库存是跨境生意里最大的资金池，也是最容易失控的地方。
本节按「预测、安全库存、补货、批量、度量、清理」的顺序展开，
这是库存决策链的天然顺序：预测错了，后面所有参数都会错。

\subsection{需求预测：从移动平均到机器学习}

\textbf{移动平均}用最近 $k$ 期的均值作为下一期预测：

\begin{equation}
\hat{y}_{t+1} = \frac{1}{k}\sum_{i=0}^{k-1} y_{t-i} .
\label{eq:ma}
\end{equation}

它的优点是几乎不需要参数、异常值影响可控；缺点是\emph{对趋势有系统性滞后}，
且对所有历史点等权，无法反映「近期更重要」。

\textbf{指数平滑}用衰减权重解决等权问题：

\begin{equation}
\hat{y}_{t+1} = \alpha y_t + (1-\alpha)\hat{y}_t, \qquad \alpha \in (0,1),
\label{eq:ses}
\end{equation}

$\alpha$ 越大越贴近近期波动。它的等价形式是「对历史做几何加权平均」，
有效记忆长度约为 $1/\alpha$ 期。$\alpha$ 的应用经验：
需求稳定取 0.1--0.2，波动大取 0.3--0.5。

\textbf{Holt-Winters} 在平滑的基础上显式建模趋势与季节：

\begin{equation}
\begin{aligned}
\ell_t &= \alpha (y_t - s_{t-m}) + (1-\alpha)(\ell_{t-1} + b_{t-1}), \\
b_t &= \beta(\ell_t - \ell_{t-1}) + (1-\beta) b_{t-1}, \\
s_t &= \gamma (y_t - \ell_t) + (1-\gamma) s_{t-m}, \\
\hat{y}_{t+h} &= (\ell_t + h\,b_t)\, s_{t-m+h},
\end{aligned}
\label{eq:holtwinters}
\end{equation}

其中 $\ell_t$ 是水平项，$b_t$ 是趋势项，$s_t$ 是周期为 $m$ 的季节项（各季节因子之和为 $m$）。
跨境场景中 $m = 12$（月度）或 $m = 52$（周度）。
Holt-Winters 是季节性商品（节庆、换季）默认的基线模型，
它最大的价值不是预测精度，而是把\emph{趋势}与\emph{季节}拆开——
这样你才能判断「这个月卖得好」是趋势真的起来了，还是只是季节性到了。
方法学细节与诊断手段可参考 \cite{hyndman2021forecasting}。

\textbf{机器学习预测}的收益与代价同样明确。可用的模型从梯度提升树（LightGBM、XGBoost）
到专用时间序列模型（Prophet、N-BEATS、DeepAR 等）。
机器学习的相对优势来自\emph{额外特征}：促销日历、价格、广告投入、
平台流量指数、竞品动作、天气与节假日。
如果拿不到这些特征，用复杂模型预测单条 SKU 的销量通常还不如 Holt-Winters 稳定可靠——
\textbf{这不是算法问题，而是样本量与信噪比问题}：
单 SKU 的历史序列常常只有几十个点，而特征维度可能上百。

\begin{table}[htbp]
\centering
\caption{四类预测方法的适用边界（选型时应按 SKU 的销量形态分组，而不是全站统一）}
\label{tab:sup-forecast}
\small
\begin{tabularx}{\textwidth}{L{2.2cm} C{1.6cm} C{1.8cm} C{1.7cm} Y}
\toprule
\tabhead{方法} & \tabhead{数据要求} & \tabhead{可解释性} & \tabhead{异常鲁棒性} & \tabhead{适用场景} \\
\midrule
移动平均 & 低（$\ge 2k$ 期） & 高 & 中 & 稳定长尾 SKU、快速建立基线 \\
指数平滑 & 低 & 高 & 低（需调优 $\alpha$） & 无季节、波动平缓的日用消耗品 \\
Holt-Winters & 中（$\ge 2$ 个完整周期） & 中 & 低 & 强季节性、趋势明显的品类 \\
机器学习加特征 & 高（多 SKU 联合训练） & 低 & 高 & SKU 数量大、促销频繁、有广告与流量特征 \\
\midrule
\multicolumn{5}{@{}>{\raggedright\arraybackslash}p{\linewidth}@{}}{\textbf{工程实践：按 ABC-XYZ 分类分策略}——
A 类高价值 SKU 用复杂模型加人工复核，C 类长尾用简单规则，间歇需求（大量零销量）用 Croston 类方法}\\
\bottomrule
\end{tabularx}
\end{table}

\subsection{误差度量：MAPE 的陷阱与 bias 的价值}

预测做完了必须度量误差，否则无法判断模型是好是坏，也无法把误差反馈进安全库存。

\begin{equation}
\text{MAPE} = \frac{1}{n}\sum_{t=1}^{n}\frac{\abs{y_t - \hat{y}_t}}{\abs{y_t}},
\qquad
\text{bias} = \frac{1}{n}\sum_{t=1}^{n}(y_t - \hat{y}_t),
\qquad
\text{WMAPE} = \frac{\sum_t \abs{y_t - \hat{y}_t}}{\sum_t y_t} .
\label{eq:mapebias}
\end{equation}

\textbf{MAPE 的陷阱在低销量 SKU 上极其致命}：
真实销量 2 件、预测 1 件，MAPE 是 50\%；
真实销量 2 件、预测 4 件，MAPE 是 100\%。
而对一个日销千件的 SKU 来说，误差 20\% 都属于不可接受。
把两类 SKU 放在同一个 MAPE 分母里平均，指标会完全失去意义。
因此工程上普遍用\emph{聚合层面的 WMAPE}（用总量做分母）
或\emph{按销量分层的 MAPE}，并对间歇性需求使用专门指标。

\textbf{bias（带符号的平均误差）度量的东西与 MAPE 完全不同}：
MAPE 度量\emph{波动}，bias 度量\emph{系统性偏移}。
一个 MAPE 只有 15\% 但 bias 恒为 +8\% 的模型，意味着每个月都少订 8\% 的货，
三个月后就是持续缺货与排名下滑。
因此监控面板上必须同时有这两个指标：
\textbf{MAPE 决定安全库存要留多少，bias 决定预测参数要不要修正}。
bias 的绝对量超过 MAPE 的 25\% 时（经验阈值），
应当优先排查是否存在漏掉的结构性因素（促销、竞品退出、平台流量政策变化），
而不是去调平滑系数。

\subsection{安全库存与服务水平：$z$ 值的经济学}

\term{安全库存}{Safety Stock} 是为了吸收\emph{需求波动}与\emph{提前期波动}而额外持有的库存。
标准形式（同时考虑需求与提前期不确定性）是：

\begin{equation}
SS = z \cdot \sqrt{LT \cdot \sigma_d^2 + d^2 \cdot \sigma_{LT}^2},
\label{eq:ss}
\end{equation}

其中 $LT$ 是平均提前期，$\sigma_d$ 是单位时间需求的标准差，
$d$ 是单位时间平均需求，$\sigma_{LT}$ 是提前期的标准差，
$z$ 是与目标服务水平对应的正态分位数。
这个式子最常见的三个算错方式：

\textbf{(1) 把 $\sigma_d$ 算成了「整体需求的标准差」，而不是「单期需求的标准差」。}
$\sigma_d$ 必须是\emph{单期}需求的标准差，再由 $\sqrt{LT}$ 放大。
如果直接用月度总需求的标准差代入，得到的安全库存会大出数倍。

\textbf{(2) 完全忽略了 $\sigma_{LT}$ 项。}
跨境场景里提前期的波动往往比需求波动更大：
海运 25--45 天、清关 3--15 天、入仓 2--20 天（旺季排队），
$d^2\sigma_{LT}^2$ 这一项经常占总方差的一半以上。
忽略它等于假设供应链永远准时，这是最昂贵的乐观。

\textbf{(3) $z$ 值与服务水平的口径不一致}。
「服务水平 95\%」至少有三种含义：\emph{周期服务水平}
（一个补货周期内不缺货的概率）、\emph{满足率}
（需求被即时满足的比例）、\emph{订单满足率}。
只有第一种直接对应 $z$，后两种需要额外的换算。
口径不统一是「明明设了 95\%，为什么还是天天断货」的常见答案。

$z$ 值来自正态分布的分位数，取几个常用值：
90\% 对应 $z = 1.282$，95\% 对应 $z = 1.645$，
97.5\% 对应 $z = 1.960$，99\% 对应 $z = 2.326$。

\begin{equation}
z = \Phi^{-1}(SL),
\qquad
\min_{SL}\ \Big[\, \underbrace{h\,\sigma_{LTd}\,\Phi^{-1}(SL)}_{\text{持有成本}} \;+\; \underbrace{c_u\, D\,(1-SL)}_{\text{缺货成本}} \,\Big],
\label{eq:zservice}
\end{equation}

其中 $h$ 是单位持有成本率，$\sigma_{LTd}$ 是提前期内需求的标准差，
$c_u$ 是单位缺货成本（含损失的毛利、排名下滑与客户终身价值损失），$D$ 是年需求。
这个最优化问题的解给出了那个反直觉结论的来源：
从 95\% 提到 99\%，$z$ 从 1.645 增到 2.326，\textbf{安全库存增加 41\%}，
而缺货概率只从 5\% 降到 1\%（少了 4 个百分点）。
所以提高服务水平是否值得，\emph{完全取决于缺货成本 $c_u$ 有多大}。
对于排名敏感、恢复成本高的新品期 SKU，$c_u$ 很大，99\% 可能值得；
对于长尾、可替代性强的老品，95\% 甚至 90\% 才是经济最优。

\subsection{再订货点与 EOQ：什么时候订、订多少}

\term{再订货点}{Reorder Point, ROP} 回答「什么时候订」：

\begin{equation}
ROP = d\cdot LT + SS,
\qquad
T^{\star} = d\,(R + LT) + z\,\sigma_{R+LT} \ \ (\text{定期盘点时用目标库存水平}),
\label{eq:rop}
\end{equation}

其中 $R$ 是盘点周期。连续盘点（系统实时看库存）用 ROP；
定期盘点（每周一统一补货）用目标库存水平 $T^{\star}$。
两者的差别在于 $R$ 期内也要覆盖需求，所以定期盘点的平均库存更高——
这是「流程方便」与「资金占用」之间的直接交换。

\term{经济订货批量}{Economic Order Quantity, EOQ} 回答「订多少」：

\begin{equation}
EOQ = \sqrt{\frac{2DS}{H}} = \sqrt{\frac{2DS}{h\,C}},
\label{eq:eoq}
\end{equation}

\begin{equation}
TC(Q) = \underbrace{\frac{D}{Q}S}_{\text{订货成本}} + \underbrace{\frac{Q}{2}H}_{\text{持有成本}} + \underbrace{D\,C}_{\text{采购成本}},
\qquad
TC(EOQ) = \sqrt{2DSH} + DC .
\label{eq:tc}
\end{equation}

其中 $D$ 是年需求，$S$ 是单次订货成本（含沟通、验货、报关、头程固定费用），
$H = hC$ 是单位年持有成本，$h$ 是持有成本率（资金成本加仓储、保险与损耗，经验上 0.15--0.30）。
EOQ 的核心性质是：\textbf{在最优点上，年订货成本恰好等于年持有成本}。
这个性质极其实用——如果你算出某 SKU 的订货成本是持有成本的三倍，
说明订货批量太小、频率太高；反之则说明批量过大。

EOQ 的三个工程修正必须知道：
(1) \textbf{有数量折扣时}，EOQ 不再是最优，
要比较「每个折扣区间的最小可行批量」与「EOQ」的总成本，取最低者（这也是 \cref{eq:tierprice} 阶梯报价的意义所在）；
(2) \textbf{有 MOQ 与整箱约束时}，实际批量取 $\max(EOQ, \text{MOQ})$ 并向上对齐到整箱；
(3) \textbf{EOQ 假设需求确定}，而现实中需求是随机的，
所以真实系统里 ROP 与 EOQ 是\emph{两个独立答案的组合}：
ROP 管触发时机，EOQ 管单次数量。经典教材 \cite{silver1998inventory} 对 $(s, Q)$ 与 $(R, S)$ 策略有完整推导。

\begin{figure}[htbp]
\centering
\begin{tikzpicture}[x=0.78cm, y=0.0052cm, font=\sffamily\tiny]
\draw[-{Stealth}, SlateGray] (0,0) -- (15.6,0) node[below right, font=\sffamily\tiny] {时间 $t$};
\draw[-{Stealth}, SlateGray] (0,0) -- (0,1220) node[above, font=\sffamily\tiny] {库存量};
\draw[dashed, AmberDeep] (0,400) -- (15.2,400);
\node[AmberDeep, right, font=\sffamily\tiny] at (15.2,400) {ROP};
\draw[dashed, CrimsonDeep] (0,150) -- (15.2,150);
\node[CrimsonDeep, right, font=\sffamily\tiny] at (15.2,150) {SS};
\draw[very thick, NavyMid] (0,1000) -- (5,150) -- (5,1000) -- (10,150) -- (10,1000) -- (14.4,320);
\fill[AmberDeep] (3,400) circle (1.4pt);
\node[AmberDeep, above left, font=\sffamily\tiny] at (3,470) {触发补货};
\draw[{Stealth}-{Stealth}, SlateGray] (3,780) -- (5,780);
\node[SlateGray, above, font=\sffamily\tiny] at (4,790) {提前期 $LT$};
\draw[{Stealth}-{Stealth}, TealDeep] (4.35,150) -- (4.35,1000);
\node[TealDeep, left, font=\sffamily\tiny, align=right] at (4.28,575) {订货批量\\$Q \approx EOQ$};
\draw[{Stealth}-{Stealth}, VioletDeep] (0,1180) -- (5,1180);
\node[VioletDeep, above, font=\sffamily\tiny] at (2.5,1200) {补货周期};
\node[draw=RuleGray, fill=BoxBGAlt, rounded corners=2pt, inner sep=3pt,
      font=\sffamily\tiny, text width=11.4cm, align=left] at (7.6,-300)
  {三个参数的分工：\emph{ROP} 决定何时下单（覆盖提前期内的需求与波动），
   \emph{订货批量} 决定一次订多少（平衡订货成本与持有成本），
   \emph{SS} 决定水位下限（由服务水平与提前期不确定性共同决定）。
   任何把三者混为一谈的「经验备货量」，都无法在需求变化时解释自己为什么是这个数。};
\end{tikzpicture}
\caption{补货决策的库存水位：SS、ROP 与 EOQ 各自回答一个不同的问题}
\label{fig:sup-inventory}
\end{figure}

\subsection{周转率、售罄率与资金占用}

库存指标里最有用的三个是周转率、周转天数与售罄率：

\begin{equation}
\text{周转率} = \frac{\text{COGS}}{\overline{I}},
\qquad
\text{DIO} = \frac{365}{\text{周转率}},
\qquad
C_{\text{资金}} = \overline{I}\cdot r_{\text{资金}} ,
\label{eq:turnover}
\end{equation}

\begin{equation}
\text{售罄率} = \frac{\text{期间销量}}{\text{期初库存} + \text{期间进货}} .
\label{eq:sellthrough}
\end{equation}

其中 $\overline{I}$ 是平均库存（注意\emph{成本口径与售价口径不能混用}，
分子分母必须同口径），$r_{\text{资金}}$ 是资金成本率。
跨境实物电商的周转率经验区间差异极大：
快消与耗材类可以做到年周转 8--12 次（DIO 约 30--45 天），
家居与服饰常在年周转 2--4 次（DIO 约 90--180 天），
而一旦出现滞销，DIO 会迅速恶化到 300 天以上（示范量级）。
\emph{周转率的每一档恶化，都直接对应一笔被锁死的现金}：
假设平均库存成本口径为 500 万元，DIO 从 90 天恶化到 150 天，
在销量不变的情况下相当于多占用约 330 万元现金。

售罄率是\emph{单品级}的健康度指标，它对新品尤其重要：
新品上架 30 天的售罄率低于 20\% 通常意味着选品或定价有问题（示范量级），
而高于 80\% 则意味着可能备货不足、正在丢失销售额。
把售罄率与库龄结构放在一起看，才能区分「卖得好」与「卖不动」。

\subsection{滞销处理：什么时候该止损}

滞销处理的本质是一个\emph{止损时点}问题。持有一个滞销 SKU 的每一期，
你都在支付仓储费、资金成本与残值衰减，同时承担平台库龄附加费的风险。
决策规则可以显式写成「立即清仓」与「再持有一期」的比较：

\begin{equation}
\underbrace{(1-\delta)P_{\text{清}}}_{\text{立即清仓的净回收}}
\;\ge\;
\underbrace{\Prob(\text{售出})\bigl(P_{\text{降价}} - h_{\text{持有}}\bigr)
+ \bigl(1-\Prob(\text{售出})\bigr)(1-\delta)P_{\text{清}}
- c_{\text{库龄}}}_{\text{再持有一期的期望回收}} ,
\label{eq:clearance}
\end{equation}

其中 $\delta$ 是清仓渠道的折价率，$P_{\text{清}}$ 是清仓基准价，
$P_{\text{降价}}$ 是降价后的售价，$h_{\text{持有}}$ 是一期的持有成本，
$c_{\text{库龄}}$ 是该期的库龄附加费。
把式子整理后得到一个很有操作性的直觉：
\textbf{当「一期持有成本加库龄附加费」超过「降价能多卖出的期望毛利」时，就应该清仓。}
在实操中，团队很少真的去算这个式子，而是用\emph{库龄分级}把它固化成规则。

\begin{table}[htbp]
\centering
\caption{按库龄分级的滞销处置规则（阈值为示范量级，须按品类动销速度与平台费率校准）}
\label{tab:sup-aged}
\small
\begin{tabularx}{\textwidth}{C{1.5cm} L{2.4cm} L{3.0cm} Y}
\toprule
\tabhead{库龄} & \tabhead{状态判定} & \tabhead{建议动作} & \tabhead{要点与风险} \\
\midrule
0--30 天 & 正常新品期 & 观察售罄率与转化率，不做价格干预 & 避免过早降价伤害价格锚与后续促销空间 \\
31--60 天 & 预警 & 优化商品页与广告结构，测试小幅降价 5\%--10\% & 优先修转化率，而不是先降价 \\
61--90 天 & 滞销 & 组合捆绑、站内促销、站外清理渠道试点 & 捆绑要选\emph{能提升客单价}的搭配，不当废品处理 \\
91--180 天 & 重度滞销 & 清仓渠道整体出货，或跨市场转移到其他站点 & 计算到手净回收，避免「清仓费高于货值」 \\
180 天以上 & 尾部处置 & 捐赠、回收或销毁，接受残值损失 & 部分市场对捐赠有税务优惠，但须严格核对当地最新规则 \\
\bottomrule
\end{tabularx}
\end{table}

关于捐赠与销毁，需要特别注意合规表述：
不同国家对库存捐赠的税务处理差异很大，需要满足受赠机构资质、
估值方法与留档要求等条件，\textbf{具体规则以当地税务机关与平台的最新公告为准}，
不能凭行业传闻做账务处理。

\begin{pitfall}[安全库存算错，导致旺季后大量滞销]
\textbf{现象}：某节庆品类 SKU 在旺季前把安全库存设得很高，
旺季确实没有断货，但旺季后 3 个月仓库里还剩 40\% 的库存，
DIO 从 60 天飙到 240 天，第二年被迫清仓，全年利润被吃掉三分之一。

\textbf{根因}：
\begin{enumerate}
  \item \textbf{把促销峰值当成需求基线}。用 11--12 月的销量计算 $d$ 与 $\sigma_d$，
        这两个参数天然包含了数倍的季节性放大，直接代入 \cref{eq:ss} 会把安全库存放大约一个量级。
  \item \textbf{用整体需求方差代替提前期需求方差}。$\sigma_d$ 用月度总量算法，
        没有按提前期长度做缩放，导致 $LT\sigma_d^2$ 项偏大。
  \item \textbf{忽略 $\sigma_{LT}$ 的方向性}。旺季提前期\emph{普遍延长}而非对称波动，
        正态假设低估了「旺季尾巴」风险，团队因此用平均交期而不是最坏交期做规划。
  \item \textbf{把服务水平一刀切}。所有 SKU 都用 99\%，包括那些毛利低、可替代性强的长尾品，
        于是资金被大量锁在低价值库存上。
  \item \textbf{旺季结束没有立刻调参}。安全库存参数一旦设定就整季沿用，
        需求已经回落到基线，库存水位却还在峰值水平。
\end{enumerate}

\textbf{排查方法}：
(1) 建立「\emph{去季节化的基线需求}」作为所有库存参数的输入，
季节性因子单独维护并按月重估；
(2) 安全库存参数\emph{按 ABC-XYZ 分组}设置，A 类高价值且需求稳定的 SKU 用高服务水平，
C 类低价值长尾降到 90\%；
(3) 用提前期与预测误差的\emph{实际分布}替换理论正态假设，
直接取「历史提前期分布的第 95 分位」作为规划提前期；
(4) 旺季结束后两周内强制触发一次参数重估，
并在系统里加一条「季节性窗口外不得沿用峰值参数」的校验；
(5) 把周转天数与库龄结构放进周报，
让滞销在 60 天而不是 180 天时就被看见。
\end{pitfall}

\begin{formulacard}[库存计划的五个核心量]
\begin{itemize}
  \item 安全库存：$SS = z\sqrt{LT\sigma_d^2 + d^2\sigma_{LT}^2}$（\cref{eq:ss}）——
        同时覆盖需求波动与提前期波动，两个方差项都不可省。
  \item 再订货点：$ROP = d\,LT + SS$（\cref{eq:rop}）——
        只回答「什么时候下单」。
  \item 经济批量：$EOQ = \sqrt{2DS/H}$（\cref{eq:eoq}）——
        只回答「一次订多少」；最优点上订货成本等于持有成本。
  \item 预测误差：MAPE 量波动、bias 量偏移（\cref{eq:mapebias}）——
        两者必须同时监控，只盯一个必然踩坑。
  \item 资金占用：$C_{\text{资金}} = \overline{I}\, r_{\text{资金}}$，
        且 $DIO = 365/\text{周转率}$（\cref{eq:turnover}）——
        库存决策最终是一道现金流题，不是一道缺货题。
\end{itemize}
\end{formulacard}

\section{履约与仓配}
\label{sec:sup-fulfill}

\subsection{海外仓选址：靠近港口还是靠近消费者}

\kw{港口仓}（如洛杉矶与长滩、萨凡纳、休斯顿、纽约与新泽西）
的优势是头程干线短、入库快、单位仓储成本低；
劣势是尾程派送距离远，跨区派送的\emph{分区费率}会显著抬高最后一公里成本。
\kw{消费市场仓}（如内陆枢纽城市、区域分仓）的优势是尾程距离短、时效快、
退货处理方便；劣势是干线要先运进来（可能多一次中转）、
仓储租金与人工更贵、且需要更精细的库存分配。

这个选择可以用一张总成本表达式来形式化。设仓网节点集合为 $\mathcal{K}$：

\begin{equation}
C_{\text{网络}} = \sum_{k\in\mathcal{K}}
\Big( \underbrace{f_k}_{\text{仓储固定}} + h\,\overline{I}_k
+ \sum_{z} q_{kz}\,c_{\text{尾程}}(z) + c_{\text{干线},k}
+ \lambda\, \E[\text{缺货损失}_k] \Big),
\label{eq:warehouse}
\end{equation}

其中 $q_{kz}$ 是从仓 $k$ 发往分区 $z$ 的订单量，$c_{\text{尾程}}(z)$ 是分区费率，
$\lambda$ 是缺货惩罚的权重系数。
这个式子的实际用法是做\emph{情景对比}：
「单港口仓」「双港口仓」「港口仓加内陆仓」三种方案，
把历史订单按收货邮编聚合成分区分布，代入各自的费率表算总成本。
多数情况下结论是：\textbf{订单密度高的区域值得设区域仓，密度低的区域靠港口仓覆盖}，
而不是简单地「全部靠近消费者」。

选址还要考虑四个常被忽略的因素：
(1) \textbf{税务关联}——美国多数州以「年销售额或交易笔数达到门槛」判定经济关联，
触发后需在该州注册并代收销售税，各州门槛与规则差异很大，须以各州官方最新规定为准；
(2) \textbf{旺季拥堵}——港口仓在旺季的入库排队时间可能从 2 天变成 20 天；
(3) \textbf{退货便利}——退货包裹的寄回地址如果跨区，逆向成本会翻倍；
(4) \textbf{人工与租金}——两者在区域间的差异常常超过尾程节省的金额。

\begin{table}[htbp]
\centering
\caption{港口仓与消费市场仓的对比（成本项为结构性对比，具体费率以服务商最新报价为准）}
\label{tab:sup-warehouse}
\small
\begin{tabularx}{\textwidth}{L{2.4cm} L{4.2cm} Y}
\toprule
\tabhead{维度} & \tabhead{港口仓} & \tabhead{消费市场仓} \\
\midrule
头程干线 & 短，整柜直达，单位成本低 & 长，需二次中转或内陆运输 \\
入库时效 & 快，旺季易拥堵 & 相对平稳，但受干线时效影响 \\
仓储租金 & 低（仓储用地供给充足） & 高（核心区域租金与人工更贵） \\
尾程分区 & 跨区比例高，分区费率上浮明显 & 分区集中在 1--4 区，尾程成本低 \\
时效与转化 & 偏远地区妥投慢，影响评分与复购 & 主流市场 2--3 日达，转化率与复购更好 \\
退货处理 & 逆向距离长，处理周期长 & 逆向距离短，翻新重上架更快 \\
库存分散风险 & 单点集中，缺货与滞销同时存在 & 多节点分散，需要更精细的调拨能力 \\
适用形态 & 大件、低频、低时效敏感商品 & 高频、高毛利、强时效敏感商品 \\
\bottomrule
\end{tabularx}
\end{table}

\subsection{头程：拼柜方式与时效成本权衡}

头程方式的成本与时效差异是数量级的：
海运整柜（FCL）单位成本最低但需要凑满一个柜；
海运拼箱（LCL）灵活但单位成本更高，且多一次拆箱与分拨；
空运与快递最快但单位成本可能是海运的数倍；
中欧班列与卡航介于两者之间。
\textbf{所有具体运价随油价、旺季、航线与船公司政策剧烈波动，
必须以货代的最新报价为准}，本节只讨论决策逻辑。

一个 40 尺高柜的可用装载容积通常在 60--68 CBM 之间（示范量级，实际取决于包装与堆码效率），
20 尺柜约 25--28 CBM。判断「要不要拼满一柜」的临界点在于：
多出的库存持有成本是否低于整柜带来的单位运费节省。

空运补货是否划算，可以用一个显式的差额式子判断：

\begin{equation}
\Delta = \underbrace{(c_{\text{空}} - c_{\text{海}})}_{\text{单位运费差额}}
- \underbrace{h\cdot\frac{\Delta T}{365}\cdot C}_{\text{少占用的资金成本}}
- \underbrace{\Prob(\text{缺货})\cdot \pi_{\text{单位}}\cdot \frac{\Delta T}{T_{\text{销售窗}}}}_{\text{少损失的销售毛利}} ,
\label{eq:freight}
\end{equation}

其中 $\Delta T$ 是空运比海运节省的天数，$C$ 是单位成本，$T_{\text{销售窗}}$ 是销售窗口长度。
\textbf{只有当 $\Delta < 0$ 时空运才划算。}
把典型数字代进去（示范量级）：单位运费差额 6 美元/kg、
商品 0.5 kg 即差额 3 美元；单位成本 8 美元、节省 25 天、资金成本率 20\%，
则资金节省约 0.11 美元；缺货概率 30\%、单位毛利 8 美元、销售窗 90 天，
则少损失的销售约 0.67 美元。合计节省不到 0.8 美元，远小于 3 美元的运费差额——
\textbf{对低货值、高重量商品，空运补货几乎永远不划算}；
反过来，高货值（如单价 300 美元以上的 3C）、强时效（如节庆窗口）、
或已经断货正在丢失排名的情况下，空运的账才算得过来。
这就是「用数据决定空运，而不是用焦虑决定空运」。

\subsection{尾程与最后一公里}

尾程派送是履约成本中占比最高的环节，公开的物流行业讨论中，
「最后一公里」通常占整体物流成本的 40\%--55\%（示范量级，因商品与市场而异）。
理解尾程成本需要掌握四个概念：

\textbf{(1) 分区（Zone）}。承运商按起运地与目的地的距离划分区域，
分区越大费率越高，且不是线性增长而是阶梯式跃升。
这就是 \cref{eq:warehouse} 里 $c_{\text{尾程}}(z)$ 的由来，
也是区域仓能省钱的数学原因。

\textbf{(2) 计费重（DIM Weight）}。承运商按「实重与体积重取大者」计费，
体积重系数因服务类型而异（常见的做法是长宽高之积除以 139 得到磅数）。
这一点与选品直接相关：\textbf{轻抛货在尾程上的成本劣势最大}。

\textbf{(3) 附加费结构}。住宅派送附加费、燃油附加费、
旺季附加费、超尺寸与超重附加费、地址更正费、以及退货标签费。
这些附加费合计可以占到基础运费的 15\%--40\%（示范量级），
做成本测算时\emph{必须按最坏情况把附加费算进去}，否则旺季会被打穿。

\textbf{(4) 时效与转化的关系}。时效不只是成本问题。
妥投时效从 5 天缩短到 2 天，对转化率与复购率的影响在多个公开行业报告中被反复讨论，
量级上常常相当于\emph{几个百分点的转化率提升}。
因此尾程决策的正确形式不是「选最便宜的承运商」，而是在
\emph{成本与时效}的曲线上选择使「履约成本加时效造成的转化损失」最小的那一点。

\subsection{退货逆向物流与 Returnless Refund}

退货率是跨境财务模型里最隐形的杀手，因为它是\emph{乘性}的。
退回来的商品还要经历一次完整的成本链：逆向运输、清关（跨境退货时）、
检测、翻新、重新包装、重新入仓，而残值往往大幅折损。

\begin{table}[htbp]
\centering
\caption{退货率基准与处置方式（量级为公开行业讨论的常见区间，须以自有数据为准）}
\label{tab:sup-return}
\small
\begin{tabularx}{\textwidth}{L{2.4cm} C{2.2cm} L{3.4cm} Y}
\toprule
\tabhead{类目} & \tabhead{退货率量级} & \tabhead{典型退货原因} & \tabhead{处置要点} \\
\midrule
服饰鞋帽 & 20\%--40\% & 尺码不符、色差、面料手感 & 尺码表与实拍图是最高杠杆的降退货手段 \\
3C 与配件 & 5\%--10\% & 功能故障、兼容性、描述不符 & 出厂可靠性抽测加兼容性清单 \\
家居家具 & 5\%--15\% & 破损、尺寸不符、色差 & 包装抗跌测试，大件优先走仅退款 \\
美妆个护 & 3\%--8\% & 过敏、气味、漏液 & 卫生原因通常不可二次销售，残值近零 \\
耗材与配件 & 2\%--6\% & 不适配、重复购买 & 退货成本低，但要注意低价值商品的逆向不经济 \\
\bottomrule
\end{tabularx}
\end{table}

\textbf{Returnless Refund}（仅退款不退货）是针对低价值商品的经典优化：
当退回商品的处理成本高于它的残值回收时，让客户保留商品反而更划算。
判定的临界条件可以写得很干净：

\begin{equation}
\text{选择仅退款} \iff
\underbrace{c_{\text{逆向}} + c_{\text{检测翻新}}}_{\text{退回的处理成本}}
>
\underbrace{\rho\, V_{\text{残值}}}_{\text{退回后可回收的价值}} ,
\label{eq:return}
\end{equation}

其中 $\rho$ 是残值回收率（受卫生、翻新成本、二次销售概率影响），$V$ 是商品成本。
代入示范量级：一件成本 12 美元的商品，逆向运费 6 美元、检测翻新 2 美元，
若 $\rho = 0.5$ 则回收价值 6 美元，处理成本 8 美元大于 6 美元，应选择仅退款；
若 $\rho = 0.9$（未拆封可原样重售）则回收 10.8 美元大于 8 美元，应要求退回。
\emph{这个规则必须按 SKU 分别设定阈值，并在系统里自动执行}，
否则客服会在每个退货请求上做一次情绪化的判断。

\begin{casebox}[案例：家居大件在平台仓与第三方海外仓之间的成本对比（示范量级）]
\textbf{背景}：一款 5 kg、包装尺寸约 $60\times40\times30$ cm 的家居收纳类商品，
售价 79.99 美元，目标是美国市场。团队在「平台仓（FBA 类履约）」与「第三方海外仓加自配送」之间做选择。

\textbf{成本结构对比（每月 3000 单，示范量级）}：
\begin{itemize}
  \item \textbf{履约费}：平台仓按尺寸与重量分段计费，该尺寸段单件约 9--13 美元；
        第三方仓拣配打包约 4--6 美元，但需加上自配送的面单成本（分区 1--4 约 8--12 美元），
        合计约 12--18 美元。这一项上平台仓通常\emph{更便宜}。
  \item \textbf{仓储费}：平台仓按体积计月费，旺季（10--12 月）费率显著高于淡季，
        并有 181 天以上的库龄附加费；第三方仓约 0.5--0.8 美元每立方英尺每月的量级。
        该商品体积约 4.2 立方英尺，淡季差异不大，但\emph{旺季与库龄附加费会拉开差距}。
  \item \textbf{流量与转化}：平台仓商品享有配送时效标签与会员流量倾斜，
        公开行业讨论中，这对转化率的影响常常是几个百分点量级，
        按 79.99 美元客单算，相当于每单数美元的隐性收益。
  \item \textbf{退货与逆向}：平台仓提供标准退货处理，但会收取退货处理费；
        第三方仓可以自行决定仅退款阈值，灵活性更高。
  \item \textbf{库存灵活性}：第三方仓在滞销时可直接转售或调拨到其他渠道，
        平台仓则需支付移除费或弃置费。
\end{itemize}

\textbf{结论形态}：该品类的正确解多半是\emph{混合}——
主销 SKU 放平台仓吃流量与时效，长尾与季节性 SKU 放第三方仓保留灵活性，
同时用第三方仓承接平台仓的移仓需求以规避移除费与长期仓储费。
\textbf{真正的决定因素不是单件履约费，而是「流量收益加时效收益」与「仓储与库龄风险」之间的权衡。}

\textbf{可迁移的结论}：在做履约方案对比时，必须把\emph{滞销情景}一起算进去。
只算「卖得好」的情景，所有方案都会指向平台仓；
一旦把「卖不动」情景下的库龄附加费、移除费与残值损失计入，
混合方案的期望成本通常更低。具体费率、分段与附加费规则每年都在调整，
\textbf{务必以平台与服务商的最新费率表为准}。
\end{casebox}

\subsection{平台仓的费用结构：哪些是「看着便宜、实际很贵」的}

平台仓（以主流电商平台的履约服务为例）的费用通常由五部分构成，
理解这个结构比记住具体数字重要得多：

\begin{table}[htbp]
\centering
\caption{平台仓费用的五个组成部分与治理动作（费率结构与具体金额以平台官方最新公告为准）}
\label{tab:sup-fba}
\small
\begin{tabularx}{\textwidth}{L{2.6cm} L{3.4cm} Y}
\toprule
\tabhead{费用项} & \tabhead{计费基础} & \tabhead{治理动作} \\
\midrule
配送费 & 按尺寸分段与重量分段，取实重与体积重的较大者 & 优化包装尺寸（缩小 1 cm 就可能跨段），减重，合箱发货 \\
月度仓储费 & 按每日平均占用体积计月费，旺季费率显著高于淡季 & 控制入库节奏，避免在旺季前把大量货压进仓 \\
长期仓储与库龄附加费 & 按库龄分段，超过 180 至 365 天后按件或按体积加收 & 提前 30 天做库龄预警，主动清仓或移除 \\
移除与弃置费 & 按件计费，退回到指定地址或销毁 & 滞销品批量移除到第三方仓，再走多渠道清理 \\
入库配置服务费 & 按入库货件的分配方式计费，分仓越分散费用结构越复杂 & 用「合仓加少量分仓」的组合换取最低总成本，注意分仓会抬高头程 \\
\midrule
\multicolumn{3}{@{}>{\raggedright\arraybackslash}p{\linewidth}@{}}{\textbf{隐藏项：低库存水平费（库存低于安全水位时按销量加收）、退货处理费、超尺寸附加费、以及旺季附加费}}\\
\bottomrule
\end{tabularx}
\end{table}

这张表里最值得警惕的是两个「反向激励」：
\emph{低库存水平费}惩罚「库存太少」，\emph{长期仓储费}惩罚「库存太多」。
平台的设计意图是让你保持「刚好够卖」的水位，
而这正是 \cref{eq:ss,eq:rop} 要解决的问题——
库存管理不再是「多备一点以防万一」的直觉，而是一个被费率结构强制约束的优化问题。

\section{供应链韧性}
\label{sec:sup-resilience}

\subsection{牛鞭效应：为什么需求波动会被放大}

\cite{lee1997bullwhip} 指出，供应链上游观察到的需求波动总是大于下游实际的需求波动，
并归纳出四个成因：

\textbf{(1) 需求信号处理}。上游只能用\emph{订单}推断需求，而订单是你对需求的\emph{反应}，
不是你观察到的需求。当需求小幅上升时，理性做法是提高订货量并上调安全库存，
而多个环节同时这样做，放大就发生了。

\textbf{(2) 短缺博弈}。当供给紧张时，下游会\emph{超量下单}（夸大需求以争取配额）。
一旦供给恢复，这些虚增的订单瞬间消失，上游立刻从缺货转为过剩。
这是供应链里最剧烈的波动来源，也是「一次性把安全库存拉满」最危险的地方。

\textbf{(3) 批量订货}。为了摊薄订货成本与运输成本，下游倾向于整批下单，
把连续的需求变成脉冲式的订单流。

\textbf{(4) 价格波动}。促销、返利与「囤货折扣」会让订单在时间上发生转移，
产生与真实需求无关的尖峰。

在经典的「移动平均预测加订货至水平」策略下，
上游订单量的方差与终端需求方差之比可以解析地写成：

\begin{equation}
\frac{\Var(q)}{\Var(d)} = 1 + \frac{2L}{T} + \frac{2L^2}{T^2},
\label{eq:bullwhip}
\end{equation}

其中 $L$ 是提前期，$T$ 是订货周期（这是 Chen 等在 2000 年给出的经典结果，
被广泛用作牛鞭效应的量化基准）。这个式子给出了三条非常具体的治理方向：

\textbf{缩短 $L$}。放大系数对 $L$ 是二次的：提前期从 40 天压到 20 天，
放大效应可以降到原来的四分之一量级。这解释了为什么「在靠近市场的海外仓备货」
比「在国内接单再发」的波动小得多——不是因为需求变了，而是因为反应更快了。

\textbf{拉长 $T$ 也要谨慎}。从公式看，拉长 $T$ 确实降低放大系数，
但 $T$ 变大意味着平均库存更高。这是一个典型的取舍，
必须放在 \cref{eq:tc} 的总成本框架里一起优化，不能只看放大系数。

\textbf{共享信息}。公式的前提是上游只能看到订单。
如果上下游共享终端销售数据（销售点数据、平台销量、库存水位），
上游可以直接对真实需求做预测，理论放大系数可以降到 1。
这是供应商管理库存与协同计划预测补货的价值来源。
在跨境场景里，这个动作对应的具体工程是：
\emph{把你的平台销量与库存数据接口开放给核心供应商，或至少每周同步一次滚动预测}。

\begin{insight}[牛鞭效应不是「供应商不靠谱」，而是信息结构问题]
很多团队把旺季缺货与旺季后滞销归因为「供应商产能不行」或「供应商压货」。
但如果把 \cref{eq:bullwhip} 拿出来看，会发现同一批供应商在信息透明时的表现完全不同：
当你能提供\emph{去季节化的滚动预测}、明确的\emph{促销日历}与\emph{真实销量回传}时，
上游的排产决策质量会显著提升。

反过来说，\textbf{你自己对上游做的事情，你也在对下游做}。
当你在平台仓把安全库存设得过高、用「防止断货」的名义把订单提前三个月下给工厂，
你就是牛鞭效应里那个放大波动的环节，
而代价最终会以库龄附加费与清仓损失的形式回到你身上。

可迁移的判断是：\emph{凡是出现「上游忽冷忽热」的现象，先查信息流，而不是先换供应商。}
\end{insight}

\subsection{约束理论：找到真正卡住产出的那一步}

\cite{goldratt2004goal} 的\term{约束理论}{Theory of Constraints, TOC}
给出了一条适用于任何系统的改进顺序：
(1) 识别瓶颈；(2) 挖尽瓶颈（不让它闲置、不把次品喂给它）；
(3) 迁就瓶颈（其余环节按瓶颈的节奏运行）；
(4) 提升瓶颈（投资扩产或改流程）；(5) 回到第一步，因为瓶颈会转移。

把 TOC 用在跨境供应链上，需要先接受一个反直觉的事实：
\textbf{系统的有效产出由瓶颈决定，优化非瓶颈不会增加任何产出}：

\begin{equation}
\text{系统产出} = \min_{i\in\mathcal{S}}
\left(\frac{\text{可用产能}_i}{\text{单位产品占用}_i}\right),
\qquad
T_{\text{有效产出}} = P_{\text{售价}} - \text{完全变动成本},
\label{eq:toc}
\end{equation}

\emph{有效产出}的概念与会计利润的区别在于：它不摊销固定费用，
因此可以避免「为了摊固定成本而多生产」这种典型错误。

跨境场景下，瓶颈经常不在工厂，而在下面这些位置：
\textbf{验货人力}（旺季所有订单挤在同一周出货，验货排不上队）、
\textbf{头程舱位}（旺季爆舱，订舱比生产更卡）、
\textbf{海外仓入库}（旺季排队 20 天）、
\textbf{平台的入库配额}（平台仓有库存上限，超了就发不进去）、
\textbf{资金}（采购付款与平台回款之间的时间差）。
识别方法很朴素：\emph{看哪一个环节的队列最长、等待时间最久、且它的空闲代价最高}。

TOC 的实操价值在于它把改进顺序讲清楚了：
在瓶颈上投入 1 元，产出直接增加；在非瓶颈上投入 1 元，只会增加库存与等待。
很多团队在非瓶颈上砸钱（例如为了让工厂多产而增加订单），
结果只是把库存堆得更高。

\subsection{单一来源风险与产地多元化}

\kw{单一来源风险}的经典教训是：真正的风险不是「涨价」，
而是「突然不能供货」。2020--2022 年的海运拥堵与随后的库存双杀，
以及多次区域性的限电与停工，都反复演示了这一点。
评估单一来源风险至少要问四个问题：
这个物料或工序有没有可替代的供应商（含境外）、
切换需要多久（打样、认证、模具、产能爬坡）、
切换成本是多少、以及现有库存能撑多久（覆盖天数）。

产地多元化是主要对策，但它远不是「换个国家下单」这么简单。

\begin{table}[htbp]
\centering
\caption{主要产地的结构性对比（示意性总结，关税与政策以官方最新公告为准）}
\label{tab:sup-sourcing}
\small
\begin{tabularx}{\textwidth}{L{1.8cm} L{3.2cm} L{2.9cm} Y}
\toprule
\tabhead{产地} & \tabhead{相对优势} & \tabhead{相对劣势} & \tabhead{适用品类} \\
\midrule
中国 & 产业带完整度最高、打样快、MOQ 弹性大、配套齐全 & 部分品类对美关税较高、地缘政治不确定性 & 电子、家居、五金、纺织等绝大多数品类 \\
越南 & 人工成本较低、部分贸易协定下关税优惠、产能持续扩张 & 上游面料与零部件仍多依赖进口、旺季产能紧张、交期波动大 & 纺织服装、鞋类、部分家具与电子组装 \\
印度 & 劳动力充足、内需市场大、部分品类有传统工艺优势 & 基础设施与物流效率差异大、交期一致性弱、部分原料依赖进口 & 纺织、手工艺品、部分化工与医药中间体 \\
墨西哥 & 靠近美国市场、陆运 2--7 天可达、部分协定下关税优惠 & 综合成本高于亚洲、供应链配套深度不足、产能有限 & 大件家居、汽车零部件、时效敏感的中小件 \\
\bottomrule
\end{tabularx}
\end{table}

产地多元化里最容易踩的法律红线是\kw{原产地规则}。
规避关税的常见错误做法是「在中国生产、运到第三国简单换包装或贴标后出口」，
这种做法在多数贸易协定下\emph{不构成实质性改变}，
一旦被认定构成规避（反规避调查），后果是追溯补税加罚款。
\textbf{合法的多元化必须包含实质性的加工或增值环节}，
并且要保留完整的生产记录、物料清单与成本核算以备核查。
具体规则（实质性改变标准、区域价值成分比例等）必须以目的国海关与贸易协定的最新规定为准。

\subsection{关税冲击与地缘政治的应对框架}

当关税上调时，常见的第一反应是「涨价」。
但涨价能不能传导，取决于需求的价格弹性：

\begin{equation}
Q^{\star} = Q_0\left(\frac{P_1}{P_0}\right)^{-\varepsilon},
\qquad
\Delta\pi = (P_1 - c - \Delta\tau\,V_{\text{申报}})\,Q^{\star} - (P_0 - c)\,Q_0 ,
\label{eq:tariff}
\end{equation}

其中 $\varepsilon$ 是价格弹性，$V_{\text{申报}}$ 是申报价值（关税的计税基础），
$\Delta\tau$ 是税率变化。这个式子最重要的用途不是算精确数字，而是排出一个\emph{动作优先级}：

\textbf{(1) 先做成本工程，再谈涨价。}
包装减重、结构简化、材质替代、SKU 合并（把三个 SKU 变成两个）都能直接降低关税基数与运费。
经验上，一轮认真的成本工程能拿回 3--8 个百分点的成本（示范量级），
而涨价 3\% 就可能带来两位数的销量下滑。

\textbf{(2) 再考虑品类结构调整。}
把所有 SKU 按「单位重量价值」与「关税率」两个维度画散点，
优先砍掉「低单位价值加高关税」的组合。
这本质上是一次被动的选品重构，而它往往能让整体毛利率不降反升。

\textbf{(3) 然后是产地与条款调整。}
包括产地多元化的真实落地、贸易术语的选择
（\code{DDP} 与 \code{DDU} 在关税承担方与申报主体上的差异会直接影响现金流与合规责任，
术语定义见 \cite{incoterms2020}），
以及利用自贸区或保税安排的递延效应（须评估合规成本与政策风险）。

\textbf{(4) 最后才是价格传导。}
涨价要分层执行：对弹性低的忠诚客户与差异化 SKU 先涨，
对完全同质化的 SKU 用「缩量不缩价」或「减配降本」的方式过渡。

此外还要建立\kw{情景规划}：把关税上调 10\%、25\%、40\% 三种情景下的
单位经济与年度利润各算一遍，明确每种情景下的触发动作与执行负责人。
\textbf{真正致命的不是关税本身，而是关税落地后才开始讨论对策}——
供应链的调整周期以季度计，任何临时决定都只能用最贵的方式（空运、清仓、违约）来弥补。

\section{数字化：ERP、OMS、WMS 与数据看板}
\label{sec:sup-digital}

\subsection{系统选型：先定流程，再选系统}

三类系统的职责边界必须分清，否则会出现「用 ERP 管库位」这种典型的错配：

\textbf{ERP（企业资源计划）}：财务与业务的主数据底座，
负责采购、库存账、成本核算、应付应收、多币种与多主体核算。

\textbf{OMS（订单管理系统）}：订单的路由与履约编排，
负责多渠道订单接入、拆单合单、库存分配（预占）、发货与物流跟踪。

\textbf{WMS（仓储管理系统）}：仓库内部作业，
负责库位、上架、拣货、波次、复核、盘点、批次与效期管理。

选型的首要判断标准不是功能清单，而是四个约束：
\emph{SKU 量级与增长速度}、\emph{渠道数量与类型}、
\emph{是否支持组合装与物料清单拆解}、
\emph{API 开放程度与二次开发能力}。
一个常见误区是追求「一套系统管全部」，
结果是 ERP 厂商的 WMS 模块在库位管理与波次拣货上完全不够用，
最后不得不在半年内再做一次迁移。

\begin{table}[htbp]
\centering
\caption{数字化系统的选型矩阵（成本区间为示范量级，须以实际报价与实施范围为准）}
\label{tab:sup-system}
\small
\begin{tabularx}{\textwidth}{L{1.9cm} L{3.1cm} L{2.8cm} Y}
\toprule
\tabhead{系统} & \tabhead{核心能力} & \tabhead{选型关键} & \tabhead{常见失误} \\
\midrule
ERP & 多主体核算、成本核算、采购与库存账、税务 & 多币种与多主体能力、成本核算颗粒度、是否支持跨境内账 & 用表格工具过渡太久，导致账实不符且无法追溯 \\
OMS & 多渠道订单接入、库存预占、智能路由、拆合单 & 平台 API 覆盖度、限流处理、幂等与重试 & 库存未做预占，多平台并发下超卖 \\
WMS & 库位、波次拣货、批次效期、盘点、耗材管理 & 与 OMS 的库存一致性、是否支持多仓调拨 & 跳过 WMS 直接用表格管库位，规模一上来就崩 \\
BI 看板 & 周转、售罄、履约时效、缺货、退货归因 & 指标口径字典、数据延迟、下钻能力 & 每个部门各自算一套口径，开会先吵指标定义 \\
\bottomrule
\end{tabularx}
\end{table}

\subsection{与平台 API 打通：限流、增量与一致性}

与平台打通（Amazon SP-API、Shopify Admin API、TikTok Shop Partner API 等，
具体接口与配额以各平台开发者文档的最新版本为准）有三个必须处理的工程问题：

\textbf{(1) 限流与退避}。平台的 API 配额通常按「每单位时间的请求数」与「突发桶」双重限制，
错误处理必须包含指数退避与抖动，且要区分\emph{可重试错误}（限流、超时）
与\emph{不可重试错误}（参数错误、权限不足）。

\textbf{(2) 增量同步而非全量拉取}。订单与库存的同步必须基于时间戳或游标做增量，
并处理「边界重复与遗漏」：用重叠时间窗加幂等键（如订单号加版本号）来保证既不漏单也不重复扣减库存。

\textbf{(3) 库存一致性}。多渠道同时销售同一个库存池时，
一致性问题会以\emph{超卖}的形式暴露。工程上有三层防线：
第一层是\emph{预占}（下单即扣减可用量），
第二层是\emph{安全缓冲}（可用量等于实际库存减缓冲，缓冲按同步延迟与日均销量估算），
第三层是\emph{对账}（每日定时比对平台库存与实际库存，自动生成差异工单）。

\begin{lstlisting}[style=growth, language=Python, caption={补货建议生成的核心逻辑（伪代码，仅示意流程）}]
# Replenishment suggestion: ROP trigger + EOQ / MOQ rounding
def suggest_order(sku, on_hand, in_transit, frozen):
    d, sigma_d = demand_stats(sku, window=56)       # daily mean and std, deseasonalized
    lt, sigma_lt = lead_time_stats(sku, window=180)  # days, from actual receipts only
    z = z_from_service_level(sku.service_level)      # 1.645 for 95 percent
    ss = z * math.sqrt(lt * sigma_d ** 2 + (d ** 2) * sigma_lt ** 2)
    rop = d * lt + ss
    available = on_hand + in_transit - frozen
    if available > rop:
        return None                                  # within range, no action
    eoq = math.sqrt(2 * d * 365 * sku.order_cost / sku.holding_cost)
    qty = max(eoq, sku.moq)
    qty = math.ceil(qty / sku.carton_qty) * sku.carton_qty
    return Order(sku, qty, reason="below_rop", ss=ss, rop=rop)
\end{lstlisting}

这段代码里有三个细节值得单独指出：
安全库存用的是\emph{去季节化}的 $\sigma_d$；
提前期统计只使用\emph{实际收货记录}而不是供应商承诺；
批量在 EOQ 之上强制对齐 MOQ 与整箱数——
最后这一条经常被忽略，导致系统建议的批量在采购侧根本无法执行，
运营很快就放弃系统、回到手动下单，数字化投入随之作废。

\begin{lstlisting}[style=growthnoline, language=SQLx, caption={库存周转与售罄率看板的取数逻辑（示意 SQL）}]
-- Inventory turnover, days of inventory and sell-through by sku and month
WITH inv AS (
    SELECT sku_id, month_id,
           AVG(on_hand_qty * unit_cost) AS avg_inv_cost
    FROM dw.inventory_daily
    GROUP BY sku_id, month_id
),
cogs AS (
    SELECT sku_id, month_id, SUM(qty * unit_cost) AS cogs_qty
    FROM dw.sales_daily
    GROUP BY sku_id, month_id
),
flow AS (
    SELECT sku_id, month_id,
           SUM(qty_in)   AS qty_in,
           SUM(qty_sold) AS qty_sold
    FROM dw.sku_monthly_flow
    GROUP BY sku_id, month_id
)
SELECT c.sku_id, c.month_id,
       ROUND(c.cogs_qty / NULLIF(i.avg_inv_cost, 0), 2)          AS turnover,
       ROUND(365.0 * i.avg_inv_cost / NULLIF(c.cogs_qty, 0), 1)  AS dio_days,
       ROUND(f.qty_sold / NULLIF(f.qty_in, 0), 4)                AS sell_through
FROM cogs c
JOIN inv  i ON i.sku_id = c.sku_id AND i.month_id = c.month_id
JOIN flow f ON f.sku_id = c.sku_id AND f.month_id = c.month_id;
\end{lstlisting}

\subsection{数据看板与自动化补货规则}

看板的价值不在「看得见」，而在\emph{把决策规则固化下来}。
一张能驱动动作的供应链看板至少包含四组指标，且每个指标必须有明确的口径、数据源、更新频率与负责人。

\begin{table}[htbp]
\centering
\caption{供应链数据看板的指标字典（口径是该表最重要的一列）}
\label{tab:sup-dashboard}
\small
\begin{tabularx}{\textwidth}{L{2.0cm} L{4.4cm} C{1.4cm} Y}
\toprule
\tabhead{分组} & \tabhead{指标与口径} & \tabhead{频率} & \tabhead{触发动作} \\
\midrule
库存效率 & 库存周转率（成本口径 COGS 除以平均库存）、周转天数、库龄结构（0--90 天、91--180 天、181 天以上的占比） & 周 & 库龄 181 天以上占比超阈值时启动清仓流程 \\
销售健康 & 售罄率（期间销量除以期初加进货）、动销率（有销量 SKU 占比）、滞销 SKU 数与金额 & 周 & 30 天售罄率低于阈值的新品进入复审 \\
履约时效 & 订单到发货时长、发货到妥投时长、准时妥投率、取消率、入库上架时长 & 日 & 准时妥投率跌破阈值时切换承运商或增加区域仓 \\
缺货与超卖 & 缺货率、缺货损失销售额、超卖单数、库存同步延迟 & 日 & 缺货率连续超阈值自动触发紧急补货或空运评估 \\
现金与成本 & 现金转换周期、履约成本占收入比、退货率与退货净损失 & 月 & 现金转换周期恶化时优先压缩在途与库龄 \\
\midrule
\multicolumn{4}{@{}>{\raggedright\arraybackslash}p{\linewidth}@{}}{\textbf{现金转换周期：库存天数加应收天数减应付天数，它是供应链效率最终折算成的财务语言}}\\
\bottomrule
\end{tabularx}
\end{table}

\begin{equation}
\text{CCC} = \text{DIO} + \text{DSO} - \text{DPO},
\qquad
Q^{\text{建议}} = \max\Big(0,\ \big(ROP + d\cdot R\big) - \big(I_{\text{可用}} + I_{\text{在途}}\big)\Big).
\label{eq:ccc}
\end{equation}

\cref{eq:ccc} 的第二个式子是自动化补货的核心规则：
以 ROP 为触发线，以「目标水位减可用与在途」为建议量。
把这条规则工程化时，必须加上四道\emph{人工覆盖}机制，否则系统会做出荒谬决策：
(1) \textbf{冻结窗口}——大促前 14 天冻结自动补货，由人为决定；
(2) \textbf{新品保护期}——上架 60 天内不自动补货，因为此时没有任何可靠的 $\sigma_d$；
(3) \textbf{供应商约束}——产能上限、MOQ、整箱数与交期不可行的建议直接拦截；
(4) \textbf{现金流约束}——单周补货总额超过现金预算时按 ABC 优先级排序，而不是全部执行。

\begin{toolbox}[上架前 30 天的供应链检查清单]
\textbf{产品与合规}
\begin{itemize}
  \item 目标市场的强制认证是否已取得，证书编号与测试报告是否留档（以官方最新要求为准）；
  \item 包装上的标识、警示语、条码与语言是否符合目的国要求；
  \item 规格书、金样、包装稿三方签认是否完成。
\end{itemize}
\textbf{库存与计划}
\begin{itemize}
  \item 安全库存的关键参数（$d$、$\sigma_d$、$LT$、$\sigma_{LT}$、$z$）是否都有数据来源；
  \item ROP 与 EOQ 是否已录入系统，MOQ 与整箱数约束是否生效；
  \item 首批备货量的情景压力测试：需求比预期低 40\% 时，周转天数会变成多少天。
\end{itemize}
\textbf{履约与逆向}
\begin{itemize}
  \item 头程方式、入仓时间窗与旺季拥堵预案是否明确；
  \item 承运商与分区费率表是否已录入成本模型（含附加费的最坏情况）；
  \item 仅退款阈值是否按 SKU 设定并在客服系统里生效。
\end{itemize}
\textbf{数据与对账}
\begin{itemize}
  \item 订单、库存、退货三类接口的同步是否已跑通，是否有每日对账工单；
  \item 看板上的四组核心指标是否已上线，口径是否与财务一致；
  \item 是否有「测算成本与实际成本」的月度对账机制。
\end{itemize}
\textbf{退出机制}
\begin{itemize}
  \item 是否预设了「什么条件下判定失败」的阈值（如 90 天售罄率低于 25\% 且转化率低于类目均值）；
  \item 失败时的清仓、调拨与销毁路径是否已经谈好价格与时效。
\end{itemize}
\end{toolbox}

\section{本章小结与下一步}

\begin{keypoint}[本章要点回顾]
\begin{enumerate}
  \item 选品的正确形态是\kw{闸门漏斗}（\cref{fig:sup-gates}）：需求验证、竞争分析、利润测算、
        供应链可行性、打样试单，每道闸门都要有\emph{量化阈值与明确否决线}。
        否决的成本远低于试错的成本，这是本章所有方法论的共同前提。
  \item 需求验证必须\kw{四信号交叉}：站内搜索量、外部趋势、平台榜单、评论挖掘。
        其中评论挖掘的信息密度最高——高频抱怨就是差异化清单。
        竞争分析看三个数：头部评论壁垒、$CR_5$ 集中度、首页广告位占比。
  \item 利润测算的成败在\kw{全链路}（\cref{eq:unitecon}）。
        出厂价通常只占售价的三成，而履约、广告与退货合计常超过三成。
        忘掉体积重、误会关税的计税基础、低估退货残值损失，是三个最常见的致命遗漏。
  \item 打分模型的权重不能靠拍（\cref{eq:scorecard}）：要么用历史 SKU 回归，
        要么用多角色几何平均，要么用敏感性分析验证排序稳定性。
        拿不到数据来源的因子应降权或删除。
  \item 供应链可行性的四个前置量是 MOQ、起订资金、交期与定制能力。
        首单资金按 \cref{eq:moqcash} 算全，现金周期按 \cref{eq:cashcycle} 算准；
        很多「高毛利」选品死于现金周期而不是利润率。
  \item 供应商管理的核心是\kw{鉴别工厂与贸易商}、\kw{持续有活单的备份供应商}、
        以及\kw{把质量条款写进合同}。谈判时价格只是一个变量，
        账期、MOQ、交期、旺季优先排产权都是可交换的筹码。
  \item 品控要沿\kw{IQC、IPQC、OQC} 三段布防，出货按 AQL 抽样（\cref{eq:aql}）。
        \textbf{AQL 不是不良率上限，而是验收规则}：真实不良率 5\% 的批次仍可能有过半概率被放行，
        所以必须用可靠性抽测与市场端赔付条款补足抽样拦不住的长尾风险。
  \item 质量成本四分类（\cref{eq:copq}）给出了资源分配原则：
        预防 1 元可以省下内部失败约 10 元、外部失败约 100 元（经验比例）。
        \textbf{把预算前移到预防与鉴定，把外部失败当作不可接受事件}。
  \item 预测没有万能方法（\cref{eq:ma,eq:ses,eq:holtwinters}）：
        数据少用移动平均与指数平滑，有季节性用 Holt-Winters，
        有丰富特征且 SKU 量大才上机器学习。误差要同时看 MAPE 与 bias
        （\cref{eq:mapebias}）——前者决定安全库存，后者决定参数要不要修。
  \item 安全库存（\cref{eq:ss}）的三个常见错误：用整体方差代替单期方差、
        忽略提前期方差项、服务水平口径不统一。
        提高服务水平是有价格的（\cref{eq:zservice}）：
        从 95\% 到 99\%，库存涨四成而缺货只少四个百分点，
        是否值得\emph{完全取决于缺货成本的大小}，所以服务水平必须按 SKU 分层设置。
  \item ROP 管何时订、EOQ 管订多少（\cref{eq:rop,eq:eoq,eq:tc}），
        两者不可混用；EOQ 最优点上订货成本恰等于持有成本，
        而 MOQ、整箱数与阶梯价会把它推开，必须做总成本比较。
  \item 库存的最终语言是\kw{现金流}：周转率、周转天数、售罄率与资金占用（\cref{eq:turnover,eq:sellthrough}），
        再到现金转换周期 $\text{CCC} = \text{DIO} + \text{DSO} - \text{DPO}$（\cref{eq:ccc}）。
        滞销处置的关键是\emph{止损时点}（\cref{eq:clearance}），
        用库龄分级把它固化成规则，而不是每季度做一次情绪化决策。
  \item 履约是一条可以优化的成本等式：仓网总成本（\cref{eq:warehouse}）、
        头程方式权衡（\cref{eq:freight}）、尾程分区费率与最后一公里占比、
        以及退货逆向与仅退款阈值（\cref{eq:return}）。
        平台仓的费率结构内含两个反向激励（低库存水平费与长期仓储费），
        把「多备一点以防万一」的直觉彻底改造成了一个约束优化问题。
  \item 韧性来自结构而非库存：牛鞭效应（\cref{eq:bullwhip}）是信息结构问题，
        治理手段是缩短提前期、减小批量、共享终端需求信息；
        TOC（\cref{eq:toc}）告诉你只在瓶颈上投入；
        产地多元化必须包含实质性加工才合法（\cref{tab:sup-sourcing}）；
        关税冲击的正确顺序是先做成本工程、再调品类、再换产地与条款、最后才涨价（\cref{eq:tariff}）。
  \item 数字化不是买系统，而是\kw{把决策规则写进系统}：
        ERP 管账、OMS 管路由与预占、WMS 管库位与作业，
        平台接口打通要处理限流、增量与一致性，
        看板要带口径字典，自动补货（\cref{eq:ccc}）必须配四道人工覆盖机制。
        \textbf{没有人工覆盖的自动化，会在第一次大促时把库存打穿。}
\end{enumerate}
\end{keypoint}

\begin{interview}[本章对应的面试题]
\begin{enumerate}
  \item 给你一个陌生类目，请描述你从 0 到 1 的选品流程。
        你会设哪几道闸门？每道闸门的量化通过标准是什么？如果时间只够做一件事，你做什么？
  \item 一个商品纸面毛利率 70\%，上架三个月后亏损。请列出至少五个可能导致实际亏损的成本项，
        并说明你会怎么在测算阶段就把它们拦住。
  \item AQL 2.5 的含义是什么？样本量 200、$Ac=10$ 时，真实不良率 5\% 的批次被接收的概率大概是多少？
        为什么说「AQL 不是不良率上限」？你会用什么机制补足抽样检验拦不住的风险？
  \item 写出安全库存公式并逐项解释。哪个参数最容易算错？
        如果目标服务水平从 95\% 提到 99\%，库存和缺货率各变化多少？什么情况下这个提升是值得的？
  \item ROP 与 EOQ 分别回答什么问题？EOQ 最优点的性质是什么？
        当供应商给出阶梯报价时，EOQ 需要怎么修正？
  \item 什么是牛鞭效应？它的四个成因分别是什么？
        请给出至少三种缓解手段，并说明每一种手段在数学上作用于哪个变量。
  \item 海外仓应该靠近港口还是靠近消费者？给出你的决策框架。
        在什么条件下空运补货比海运划算？请写出判断式并代入一组你自己设定的数字。
  \item 关税上调 25\%，你的单位经济被侵蚀 6 个百分点。
        请按顺序给出你的应对动作，并说明为什么这个顺序不能颠倒。
\end{enumerate}
参考答案与答题框架见 \cref{app:interview}。
\end{interview}

本章讲的是\emph{把货做出来并送到用户手里}这条链路上的一切约束：资金、库存、交期、关税与合规。
它的特点与前面几章恰好相反——\textbf{反馈慢、不可逆、且几乎没有标准实验的空间}。
你不能「先备 5 万件试试看」，也很难在两周内判断一个供应链决策对错。
正因为如此，供应链的能力最终体现为\emph{用更少的库存支撑同样的销售额}，
而这恰恰是可以被度量的：周转率、售罄率、现金转换周期、缺货率。
这些指标构成了增长的另一半基础——没有它们，流量买得越多，亏损得越快。

而所有这些指标，最终都要回到一个闭环：\emph{观察数据、提出假设、做实验、沉淀结论}。
一次选品是一次实验，一次备货决策是一次实验，
一次海外仓切换也是一次实验——只是它们的周期更长、成本更高。
下一章 \cref{ch:growth} 我们就来讨论增长闭环本身：
指标体系如何搭建（从 $\GMV$ 拆解到 $\UV \times \CVR \times \AOV \times \text{复购}$）、
归因与实验设计如何在没有标准实验条件的场景下仍然做出可信判断 \cite{ellis2021growth}，
以及一个组织如何让「数据到结论」的循环转得比竞争对手更快。
如果你想先补齐跨境经营的全局框架与平台侧的规则差异，
可以先回到 \cref{ch:cbec} 与 \cref{ch:platform}；
想了解流量端的投放与内容协同，可对照 \cref{ch:meta}。
